% Options for packages loaded elsewhere
\PassOptionsToPackage{unicode}{hyperref}
\PassOptionsToPackage{hyphens}{url}
\PassOptionsToPackage{dvipsnames,svgnames,x11names}{xcolor}
\documentclass[
  11pt,
  a4paper,
]{article}
\usepackage{xcolor}
\usepackage[margin=1in]{geometry}
\usepackage{amsmath,amssymb}
\setcounter{secnumdepth}{-\maxdimen} % remove section numbering
\usepackage{iftex}
\ifPDFTeX
  \usepackage[T1]{fontenc}
  \usepackage[utf8]{inputenc}
  \usepackage{textcomp} % provide euro and other symbols
\else % if luatex or xetex
  \usepackage{unicode-math} % this also loads fontspec
  \defaultfontfeatures{Scale=MatchLowercase}
  \defaultfontfeatures[\rmfamily]{Ligatures=TeX,Scale=1}
\fi
\usepackage{lmodern}
\ifPDFTeX\else
  % xetex/luatex font selection
  \ifLuaTeX
    \usepackage{luaotfload}
    \directlua{luaotfload.add_fallback("mainfontfallback",{
      "NotoSerifBengali:mode=harf"
    })}
  \fi
  \setmainfont[,RawFeature={fallback=mainfontfallback}]{Latin Modern Roman}
  \ifLuaTeX
    \usepackage{luaotfload}
    \directlua{luaotfload.add_fallback("monofontfallback",{
      "NotoSansBengali:mode=harf"
    })}
  \fi
  \setmonofont[,RawFeature={fallback=monofontfallback}]{DejaVu Sans Mono}
\fi
% Use upquote if available, for straight quotes in verbatim environments
\IfFileExists{upquote.sty}{\usepackage{upquote}}{}
\IfFileExists{microtype.sty}{% use microtype if available
  \usepackage[]{microtype}
  \UseMicrotypeSet[protrusion]{basicmath} % disable protrusion for tt fonts
}{}
\makeatletter
\@ifundefined{KOMAClassName}{% if non-KOMA class
  \IfFileExists{parskip.sty}{%
    \usepackage{parskip}
  }{% else
    \setlength{\parindent}{0pt}
    \setlength{\parskip}{6pt plus 2pt minus 1pt}}
}{% if KOMA class
  \KOMAoptions{parskip=half}}
\makeatother
\usepackage{longtable,booktabs,array}
\usepackage{calc} % for calculating minipage widths
% Correct order of tables after \paragraph or \subparagraph
\usepackage{etoolbox}
\makeatletter
\patchcmd\longtable{\par}{\if@noskipsec\mbox{}\fi\par}{}{}
\makeatother
% Allow footnotes in longtable head/foot
\IfFileExists{footnotehyper.sty}{\usepackage{footnotehyper}}{\usepackage{footnote}}
\makesavenoteenv{longtable}
\ifLuaTeX
  \usepackage{luacolor}
  \usepackage[soul]{lua-ul}
\else
  \usepackage{soul}
\fi
\setlength{\emergencystretch}{3em} % prevent overfull lines
\providecommand{\tightlist}{%
  \setlength{\itemsep}{0pt}\setlength{\parskip}{0pt}}
% paper_preamble.tex — shared LaTeX preamble for the PCMA/SI paper set.
%
% Compiled with lualatex (NOT pdflatex) because several papers contain
% Bengali script: the CLL paper realises one cognitive state into
% English, Bengali and Banglish, so the Bengali output is EVIDENCE and
% cannot be transliterated away without destroying the point being made.
%
% Font selection is NOT done here. pandoc >= 3 has native
% `mainfontfallback` / `monofontfallback` variables that emit the
% luaotfload fallback registration in the correct position relative to
% unicode-math; doing it by hand in this header lands after unicode-math
% and breaks font loading outright. See tools/build_papers.sh.

\usepackage{amsmath}
\usepackage{amssymb}
\usepackage{mathtools}   % \xleftrightarrow etc. used in cll.md §5
\usepackage{booktabs}

% ---------------------------------------------------------------------
% Unicode symbols appearing in PROSE rather than inside $...$.
%
% The papers were written as markdown and read as markdown, so they use
% real characters — "the ω-boost", "α-child", "ΔR coherence" — instead of
% LaTeX macros. Under unicode-math those codepoints are math-mode-only
% and abort the build ("\symbballowed only in math mode"). Mapping each
% to \ensuremath keeps the markdown readable on GitHub AND compilable,
% which is the whole point of keeping one source.
% ---------------------------------------------------------------------
\usepackage{newunicodechar}

% Greek, lower case
\newunicodechar{α}{\ensuremath{\alpha}}
\newunicodechar{β}{\ensuremath{\beta}}
\newunicodechar{γ}{\ensuremath{\gamma}}
\newunicodechar{δ}{\ensuremath{\delta}}
\newunicodechar{ε}{\ensuremath{\varepsilon}}
\newunicodechar{η}{\ensuremath{\eta}}
\newunicodechar{θ}{\ensuremath{\theta}}
\newunicodechar{κ}{\ensuremath{\kappa}}
\newunicodechar{λ}{\ensuremath{\lambda}}
\newunicodechar{μ}{\ensuremath{\mu}}
\newunicodechar{π}{\ensuremath{\pi}}
\newunicodechar{ρ}{\ensuremath{\rho}}
\newunicodechar{σ}{\ensuremath{\sigma}}
\newunicodechar{τ}{\ensuremath{\tau}}
\newunicodechar{φ}{\ensuremath{\varphi}}
\newunicodechar{χ}{\ensuremath{\chi}}
\newunicodechar{ψ}{\ensuremath{\psi}}
\newunicodechar{ω}{\ensuremath{\omega}}
\newunicodechar{ν}{\ensuremath{\nu}}
\newunicodechar{ξ}{\ensuremath{\xi}}
\newunicodechar{ζ}{\ensuremath{\zeta}}
\newunicodechar{ι}{\ensuremath{\iota}}

% Greek, upper case
\newunicodechar{Δ}{\ensuremath{\Delta}}
\newunicodechar{Φ}{\ensuremath{\Phi}}
\newunicodechar{Σ}{\ensuremath{\Sigma}}
\newunicodechar{Ψ}{\ensuremath{\Psi}}
\newunicodechar{Ω}{\ensuremath{\Omega}}
\newunicodechar{Γ}{\ensuremath{\Gamma}}
\newunicodechar{Λ}{\ensuremath{\Lambda}}
\newunicodechar{Θ}{\ensuremath{\Theta}}
\newunicodechar{Π}{\ensuremath{\Pi}}

% Operators and relations
\newunicodechar{×}{\ensuremath{\times}}
\newunicodechar{·}{\ensuremath{\cdot}}
\newunicodechar{±}{\ensuremath{\pm}}
\newunicodechar{−}{\ensuremath{-}}
\newunicodechar{≈}{\ensuremath{\approx}}
\newunicodechar{≤}{\ensuremath{\leq}}
\newunicodechar{≥}{\ensuremath{\geq}}
\newunicodechar{≠}{\ensuremath{\neq}}
\newunicodechar{∈}{\ensuremath{\in}}
\newunicodechar{∂}{\ensuremath{\partial}}
\newunicodechar{∇}{\ensuremath{\nabla}}
\newunicodechar{√}{\ensuremath{\sqrt{}}}
\newunicodechar{∞}{\ensuremath{\infty}}
\newunicodechar{∝}{\ensuremath{\propto}}
\newunicodechar{∑}{\ensuremath{\sum}}
\newunicodechar{∫}{\ensuremath{\int}}
\newunicodechar{ℝ}{\ensuremath{\mathbb{R}}}
\newunicodechar{ℕ}{\ensuremath{\mathbb{N}}}
\newunicodechar{‖}{\ensuremath{\|}}
\newunicodechar{⊕}{\ensuremath{\oplus}}
\newunicodechar{⊗}{\ensuremath{\otimes}}
\newunicodechar{∀}{\ensuremath{\forall}}
\newunicodechar{∃}{\ensuremath{\exists}}
\newunicodechar{¬}{\ensuremath{\neg}}
\newunicodechar{∧}{\ensuremath{\wedge}}
\newunicodechar{∨}{\ensuremath{\vee}}

% Arrows
\newunicodechar{→}{\ensuremath{\rightarrow}}
\newunicodechar{←}{\ensuremath{\leftarrow}}
\newunicodechar{↔}{\ensuremath{\leftrightarrow}}
\newunicodechar{↑}{\ensuremath{\uparrow}}
\newunicodechar{↓}{\ensuremath{\downarrow}}
\newunicodechar{⇒}{\ensuremath{\Rightarrow}}
\newunicodechar{⇐}{\ensuremath{\Leftarrow}}
\newunicodechar{⇔}{\ensuremath{\Leftrightarrow}}

% Super/subscripts written as literal characters
\newunicodechar{²}{\ensuremath{^{2}}}
\newunicodechar{³}{\ensuremath{^{3}}}
\newunicodechar{⁴}{\ensuremath{^{4}}}
\newunicodechar{⁵}{\ensuremath{^{5}}}
\newunicodechar{⁶}{\ensuremath{^{6}}}
\newunicodechar{₀}{\ensuremath{_{0}}}
\newunicodechar{₁}{\ensuremath{_{1}}}
\newunicodechar{₂}{\ensuremath{_{2}}}
\newunicodechar{₃}{\ensuremath{_{3}}}
\newunicodechar{ᵀ}{\ensuremath{^{\mathsf{T}}}}

% Status glyphs. The papers use emoji as status markers in tables
% (implemented / failed / warning). Mapped to typographic equivalents:
% a colour emoji font in lualatex is unreliable, and the MEANING is what
% carries — a reader needs to see "this one failed", not a red cross.
\newunicodechar{✅}{\ensuremath{\checkmark}}
\newunicodechar{✓}{\ensuremath{\checkmark}}
\newunicodechar{❌}{\ensuremath{\times}}
\newunicodechar{✗}{\ensuremath{\times}}
\newunicodechar{⚠}{\textbf{!}}
\newunicodechar{🧩}{\textbf{[piece]}}
\newunicodechar{📋}{\textbf{[list]}}
\newunicodechar{🟡}{\textbf{[partial]}}
\newunicodechar{🔴}{\textbf{[dead]}}
\newunicodechar{🟢}{\textbf{[live]}}
\newunicodechar{▼}{\ensuremath{\blacktriangledown}}
\newunicodechar{►}{\ensuremath{\blacktriangleright}}
\newunicodechar{□}{\ensuremath{\square}}
\newunicodechar{█}{\rule{0.6em}{0.9ex}}
\newunicodechar{•}{\textbullet}
\newunicodechar{†}{\textdagger}
\newunicodechar{‑}{-}
\newunicodechar{‐}{-}
\newunicodechar{…}{\ldots}
\newunicodechar{।}{\textbar}   % danda: rendered by the Bengali fallback in text
\newunicodechar{⁻}{\ensuremath{^{-}}}

% Zero-width formatting codepoints. These are invisible in an editor but
% are real characters to LaTeX: an emoji written as "⚠️" is TWO
% codepoints (U+26A0 U+FE0F), so declaring only the visible one leaves a
% variation selector with no definition and the build dies on a symbol
% nobody can see. Both are mapped to nothing.
\newunicodechar{^^^^fe0f}{}    % VARIATION SELECTOR-16
\newunicodechar{^^^^200d}{}    % ZERO WIDTH JOINER
\newunicodechar{^^^^200b}{}    % ZERO WIDTH SPACE

\usepackage{longtable}
\usepackage{array}

% Comparison matrices and cost ledgers are wide. Let them shrink to the
% text block rather than overflow the margin.
\usepackage{adjustbox}

% Code blocks: wrap long lines instead of running off the page. Several
% papers quote real C++ with long identifiers.
\usepackage{fvextra}
\fvset{breaklines=true,breakanywhere=true,fontsize=\small}

\setlength{\parskip}{0.5em}
\setlength{\parindent}{0pt}

\usepackage{fancyhdr}
\pagestyle{fancy}
\fancyhf{}
\fancyfoot[C]{\thepage}
\fancyhead[L]{\small\itshape Synthetic Intelligence — research lineage}
\fancyhead[R]{\small\itshape irfan@furylogic.com}
\renewcommand{\headrulewidth}{0.2pt}
\setlength{\headheight}{14pt}
\setlength{\headsep}{18pt}
\usepackage{bookmark}
\IfFileExists{xurl.sty}{\usepackage{xurl}}{} % add URL line breaks if available
\urlstyle{same}
\hypersetup{
  pdftitle={From Artificial to Synthetic Intelligence --- Principles of Constructive Structural Cognition},
  pdfauthor={Irfan Mahir --- Independent Researcher, Dhaka, Bangladesh --- irfan@furylogic.com},
  colorlinks=true,
  linkcolor={black},
  filecolor={Maroon},
  citecolor={Blue},
  urlcolor={Blue},
  pdfcreator={LaTeX via pandoc}}

\title{From Artificial to Synthetic Intelligence --- Principles of
Constructive Structural Cognition}
\author{Irfan Mahir --- Independent Researcher, Dhaka, Bangladesh ---
irfan@furylogic.com}
\date{}

\begin{document}
\maketitle

{
\hypersetup{linkcolor=}
\setcounter{tocdepth}{2}
\tableofcontents
}
\section{From Artificial to Synthetic Intelligence --- Principles of
Constructive Structural
Cognition}\label{from-artificial-to-synthetic-intelligence-principles-of-constructive-structural-cognition}

\textbf{Irfan Mahir}\\
\emph{Independent Researcher, Dhaka, Bangladesh}\\
irfan@furylogic.com

\begin{quote}
The framework paper. Read this first; every other volume is a companion
to it.
\end{quote}

\begin{center}\rule{0.5\linewidth}{0.5pt}\end{center}

\subsection{Contents of This Volume}\label{contents-of-this-volume}

\begin{longtable}[]{@{}
  >{\raggedright\arraybackslash}p{(\linewidth - 4\tabcolsep) * \real{0.3333}}
  >{\raggedright\arraybackslash}p{(\linewidth - 4\tabcolsep) * \real{0.3333}}
  >{\raggedright\arraybackslash}p{(\linewidth - 4\tabcolsep) * \real{0.3333}}@{}}
\toprule\noalign{}
\begin{minipage}[b]{\linewidth}\raggedright
\#
\end{minipage} & \begin{minipage}[b]{\linewidth}\raggedright
Paper
\end{minipage} & \begin{minipage}[b]{\linewidth}\raggedright
File
\end{minipage} \\
\midrule\noalign{}
\endhead
\bottomrule\noalign{}
\endlastfoot
1 & From Artificial to Synthetic Intelligence: Principles of
Constructive Structural Cognition &
\texttt{synthetic-intelligence.md} \\
\end{longtable}

\begin{center}\rule{0.5\linewidth}{0.5pt}\end{center}

\newpage

\subsection{From Artificial to Synthetic Intelligence: Principles of
Constructive Structural
Cognition}\label{from-artificial-to-synthetic-intelligence-principles-of-constructive-structural-cognition-1}

\begin{quote}
\emph{``Thinking is not the selection of a token; it is the construction
of a mechanism that computes the answer.''}
\end{quote}

\begin{center}\rule{0.5\linewidth}{0.5pt}\end{center}

\subsubsection{Abstract}\label{abstract}

We propose a rigorous distinction between \textbf{Artificial
Intelligence (AI)} --- pattern matching over statistical landscapes ---
and \textbf{Synthetic Intelligence (SI)} --- active construction of
deterministic mechanisms that compute answers under structural
constraints. We argue that every existing cognitive architecture,
including our own prior work (CMA, P-CMA, HCMA, DCHL, RIFA), reduces at
small scale to some hidden form of memorization, classification, or
retrieval. We prove that this is not an implementation defect but a
consequence of a conservation law: the \textbf{Translation Capacity
Lemma}, which shows that any mapping from a continuous cognitive
substrate to a discrete semantic set of size \(N\) requires channel
capacity of at least \(\log_2 N\) bits \emph{somewhere} in the system,
and that this cost can only be relocated, never eliminated. We then
present the \textbf{Constructive Structural Cognition (CSC)} framework,
built on three axioms --- categorical-topological state representation,
non-equilibrium phase-space expansion, and active constructive reasoning
over statistical interpolation --- realized in a \textbf{tri-layer
architecture}: a physics-based substrate for coherence and identity, a
typed graph-rewriting engine for constructive derivation, and a
minimum-description-length translation table that pays only the
mandatory conservation-law tax.

Four results distinguish this version from the framework as originally
stated. \textbf{(i) Constructive Language Locking} (§5.4d): deriving
grammar from field dynamics is a category error, because grammar is a
transmission protocol rather than a physical phenomenon; moving surface
realisation into isolated deterministic language clusters yields
byte-exact output in English, Bengali and Banglish from one cognitive
state, with zero learned parameters and no possibility of drift.
\textbf{(ii) A third provenance class} (§5.4c): facts \emph{established
by experiment} against an external oracle --- the C++ compiler --- which
are neither asserted nor rule-derived, and which have corrected the
graph where mechanical extraction was wrong. \textbf{(iii) The J-A-S
cycle} (§5.4e): Einstein's induction--deduction--refutation loop as a
running mechanism, including an \textbf{impossibility register} that
distinguishes a refuted induction from a theorem left intact by a scope
escape; two of this project's own impossibility theorems were escaped
without being violated. \textbf{(iv) A mechanism-level lineage audit}
(§6.2), performed by grepping live source rather than reading papers,
which retires several components this paper previously claimed to retain
--- including the DCHL substrate on which its own headline empirical
result was obtained.

That last point is the paper's methodological claim as much as its
architectural one: a lineage accumulates claims faster than it retires
them, and the only defence is to re-derive the parts list from the code
on every revision. Sections marked \textbf{RETIRED} or
\textbf{historical} below are kept deliberately, because a negative
result that is deleted stops being evidence.

\begin{center}\rule{0.5\linewidth}{0.5pt}\end{center}

\subsubsection{1. Introduction}\label{introduction}

\paragraph{1.1 The Two Kinds of
Intelligence}\label{the-two-kinds-of-intelligence}

The last decade of AI has confused \emph{statistical fluency} with
\emph{cognition}. A transformer trained on a sufficient corpus can
produce text indistinguishable from a thinker's --- because the corpus
contains the thinker's outputs, and pattern-matching over a large-enough
distribution asymptotically reproduces them. This is Artificial
Intelligence: the reconstruction of the surface of thought from its
shadow in data.

\textbf{Synthetic Intelligence is different in kind, not in degree.} SI
does not reconstruct thought from data. SI \emph{constructs} a
mechanism, for each query, that computes the answer under the constraint
of the query. A calculator is a trivial SI system: given ``17 × 23'', it
does not sample; it derives. Coq, Lean, Isabelle are more sophisticated
SI systems: given a proposition, they construct a proof or report that
none exists. What has been missing is a general framework that lifts
this kind of active construction to natural queries, retains the
interpretability of formal systems, and does not collapse into either
symbolic AI's brittleness or statistical AI's opacity.

This paper is that framework.

\paragraph{1.2 The Three Traps of Prior
Attempts}\label{the-three-traps-of-prior-attempts}

Every prior attempt at physics-based or structural cognition ---
including our own --- collapses at small scale into one of three failure
modes:

\begin{itemize}
\tightlist
\item
  \textbf{The Classifier Trap.} Contrastive training (text-vs-random,
  coherent-vs-shuffled) trains an energy function that discriminates
  classes. At inference the state descends into the winning class's
  centroid. This is a binary classifier; nothing about ``physics-based''
  changes that.
\item
  \textbf{The Memorization Trap.} Q → A triplet training forces the
  network to reproduce specific input-output mappings. The physics is
  decoration on a lookup table.
\item
  \textbf{The Retrieval Trap.} Cosine similarity over an evolved-state
  anchor bank is attention with a different name. The evolution is a
  normaliser; the retrieval is the substance.
\end{itemize}

We show below that these are not three defects but three symptoms of the
same underlying constraint: information cannot be produced out of
nothing, and any cognitive system that reliably distinguishes \(N\)
semantic outcomes must contain \(\log_2 N\) bits of information about
which outcome to produce. Physics-based dynamics do not exempt one from
this law.

\paragraph{1.3 Contributions}\label{contributions}

\begin{itemize}
\tightlist
\item
  \textbf{The Translation Capacity Lemma} (§3), a conservation law for
  cognitive architecture.
\item
  \textbf{The Composition-Retrieval Distinction} (§3.4), an operational
  definition of ``thinker'' as \emph{composition rate over storage
  rate}, replacing the impossible target of ``no storage.''
\item
  \textbf{The Constructive Structural Cognition framework} (§4), three
  axioms distinguishing SI from AI.
\item
  \textbf{The Tri-Layer Unified Architecture} (§5), a synthesis of our
  prior lineage plus the missing constructive layer.
\item
  \textbf{Constructive Language Locking} (§5.4d, \texttt{cll.md}),
  separating cognitive freedom from linguistic protocol, with measured
  three-language output from a single triple.
\item
  \textbf{The experiment provenance class} (§5.4c), knowledge
  established by an external oracle rather than asserted or derived,
  with the cost of a repeated task falling from 3 experiments to 1.
\item
  \textbf{The J-A-S cycle and impossibility register} (§5.4e), including
  the rule that a counterexample to a THEOREM is a scope error until
  proven otherwise --- the distinction between Minsky--Papert being
  \emph{escaped} and being \emph{violated}.
\item
  \textbf{The place of prior work} (§6) as intellectual history, and
  \textbf{a mechanism-level audit} (§6.2) stating, per paper and by
  grep, what is load-bearing, what is partly salvaged, and what is dead.
  Two mechanisms (ZeroGrad, ANQ-LEA) are live while their papers are
  retired.
\item
  \textbf{Open problems and empirical program} (§7--§8), with closed
  items marked closed and abandoned items marked abandoned rather than
  pending.
\end{itemize}

\begin{center}\rule{0.5\linewidth}{0.5pt}\end{center}

\subsubsection{2. The Failure Modes of Substrate-Only
Cognition}\label{the-failure-modes-of-substrate-only-cognition}

Our prior work developed a family of \emph{physics-based cognitive
substrates}: continuous state spaces, learned potentials, Hamiltonian
dynamics, curl-free vector fields, energy-based Hebbian training. Each
substrate came with a paper claiming that its dynamics constituted
``thinking.'' We now argue --- and this is a self-critique --- that all
of them, at any scale we have empirically tested (from 33K to 1B
parameters), collapse into the three traps of §1.2.

\paragraph{2.1 The Substrate Alone Cannot
Answer}\label{the-substrate-alone-cannot-answer}

Consider the simplest instance. The DCHL potential \(\Phi(\mathbf{c})\)
is trained by decoupled contrastive Hebbian learning to have low energy
on real text windows and high energy on incoherent ones. The Hamiltonian
evolution \(\dot{\mathbf{c}} = -\nabla\Phi(\mathbf{c})\) then descends
into a nearby low-energy basin.

For any semantic question --- ``what is the capital of France?'' --- the
substrate produces an evolved state \(\mathbf{c}^*\) that is
\emph{coherent} (low \(\Phi\)) but has no discrete semantic content. It
is a point on a manifold, not an answer. Reading out ``Paris'' from
\(\mathbf{c}^*\) requires an additional map
\(\phi: \mathcal{M} \to \Sigma\) that has semantic knowledge the
substrate does not.

The Byte Energy Descent (BED) decoder {[}Mahir 2026e{]} attempted to
make this map be a discrete argmin of \(\Phi\) over the byte lattice.
Empirically, this converged to byte-mean states without
prompt-conditioned structure --- because \(\Phi\) was trained to score
whole-window coherence, not to encode prompt-answer dependency. The
retrieval-based CognitiveMemory alternative worked (it produced actual
English answers) --- but only because it stored \((Q, A)\) pairs, which
is memorization.

\paragraph{2.2 Why Every Substrate-Only Approach Reduces to One of the
Three
Traps}\label{why-every-substrate-only-approach-reduces-to-one-of-the-three-traps}

The pattern generalizes:

\begin{itemize}
\tightlist
\item
  If the substrate is trained on \emph{class-vs-class} contrast →
  classifier trap.
\item
  If the substrate is trained on \emph{specific mappings} → memorizer
  trap.
\item
  If a lookup mechanism is added on top of the substrate → retriever
  trap.
\end{itemize}

There is no fourth option that produces a \emph{specific} discrete
answer from a \emph{specific} input. This is not because the substrate
design is wrong. It is because information about which of \(N\) answers
to produce cannot arise from dynamics alone; it must be stored, in some
form, in the system's parameters or state. This is the theorem of §3.

\begin{center}\rule{0.5\linewidth}{0.5pt}\end{center}

\subsubsection{3. The Translation Capacity
Lemma}\label{the-translation-capacity-lemma}

\paragraph{3.1 Statement}\label{statement}

\textbf{Lemma (Translation Capacity).} Let \(\mathcal{M}\) be a state
space equipped with a computable dynamical system, and let \(\Sigma\) be
a discrete semantic set with \(|\Sigma| = N\). Let
\(\phi: \mathcal{M} \to \Sigma\) be a computable map that reliably
discriminates all \(N\) elements (i.e., for every \(\sigma \in \Sigma\)
there exists \(\mathbf{x} \in \mathcal{M}\) with
\(\phi(\mathbf{x}) = \sigma\), and the pre-images are separable to
within some error \(\epsilon > 0\)).

Then the total description complexity of \(\phi\) (as measured by its
Kolmogorov complexity, or equivalently by the channel capacity of any
implementation) is at least

\[
K(\phi) \geq \log_2 N - O(1).
\]

Moreover, this cost cannot be relocated \emph{out of} the system; it can
only be redistributed \emph{within} it.

\paragraph{3.2 Proof Sketch}\label{proof-sketch}

The cost is an immediate consequence of the source coding theorem: the
output \(\sigma\) carries \(\log_2 N\) bits of information, and the only
source of this information is \(\phi\) itself, since \(\mathcal{M}\) has
continuous cardinality and the dynamics are computable (thus of bounded
description). Any implementation of \(\phi\) must therefore
\emph{contain} those \(\log_2 N\) bits. Distributing them across many
parameters (as in an LLM), storing them explicitly (as in a lookup
table), or hiding them in the shape of a learned potential (as in DCHL)
does not reduce the total; it only changes the compression ratio.
\(\square\)

\paragraph{3.3 The Cost Ledger Across
Architectures}\label{the-cost-ledger-across-architectures}

The lemma allows a clean accounting of where the mandatory \(\log_2 N\)
bits sit in each architecture we have studied:

\begin{longtable}[]{@{}
  >{\raggedright\arraybackslash}p{(\linewidth - 4\tabcolsep) * \real{0.3333}}
  >{\raggedright\arraybackslash}p{(\linewidth - 4\tabcolsep) * \real{0.3333}}
  >{\raggedright\arraybackslash}p{(\linewidth - 4\tabcolsep) * \real{0.3333}}@{}}
\toprule\noalign{}
\begin{minipage}[b]{\linewidth}\raggedright
Architecture
\end{minipage} & \begin{minipage}[b]{\linewidth}\raggedright
Location of \(\log_2 N\) bits
\end{minipage} & \begin{minipage}[b]{\linewidth}\raggedright
Explicit?
\end{minipage} \\
\midrule\noalign{}
\endhead
\bottomrule\noalign{}
\endlastfoot
Transformer LM & Vocabulary embedding + softmax head & Yes \\
Frontier LLM & Distributed across \(\sim 10^{10}\)--\(10^{12}\)
parameters & No (opaque) \\
Symbolic AI & Rule table + fact base & Yes \\
CMA / P-CMA / HCMA & Shape of learned \(\Phi\) + decoder & No
(obscured) \\
DCHL + BED & Learned \(\Phi\) + argmin over bytes & Partial \\
CognitiveMemory & Anchor set + string bank & Yes \\
STAN topological readout & Winding-number → concept table & Yes \\
RIFA Rule Synthesizer & Antecedent-consequent bindings & Yes \\
\end{longtable}

\textbf{Consequence.} The right question is not \emph{``can we build a
cognitive architecture without a translation layer?''} --- the lemma
proves we cannot. The right question is: \emph{``what is the
minimum-description-length translation layer that a coherence-driven
substrate can operate, and what work can only physics-based substrate
perform that no lookup table can?''}

This reformulation is the philosophical hinge of Synthetic Intelligence.

\paragraph{3.4 The Composition--Retrieval
Distinction}\label{the-compositionretrieval-distinction}

The lemma dissolves the naïve ``storage-free thinker'' ambition but
leaves a sharper distinction intact. We now state it precisely.

\textbf{Definition (Composition rate and retrieval rate).} Let a
cognitive system \(\mathcal{S}\) have a set of stored primitives \(\Pi\)
with \(|\Pi| = P\), and let it operate on queries \(q\) by producing
derivations \(d(q)\) that are compositions of primitives of length
\(\ell(q)\).

\begin{itemize}
\tightlist
\item
  The \textbf{retrieval rate} of \(\mathcal{S}\) is
  \(R(\mathcal{S}) = |\{q : \ell(q) \leq 1\}| / |Q|\): the fraction of
  queries answered by a single primitive lookup.
\item
  The \textbf{composition rate} is
  \(C(\mathcal{S}) = \mathbb{E}_q[\ell(q)]\): the expected derivation
  length on the query distribution.
\end{itemize}

\textbf{Definition (Thinker vs.~Pattern Matcher).} A system is a
\textbf{thinker} to the extent that \(C(\mathcal{S}) \gg 1\) and the
number of \emph{reachable derivations}
\(|\{\text{compositions of length} \leq L\}|\) grows exponentially in
\(L\) while \(P\) grows only polynomially. A system is a \textbf{pattern
matcher} to the extent that \(R(\mathcal{S}) \to 1\) and
\(C(\mathcal{S}) \to 1\).

The distinction is \emph{not} about the presence or absence of storage
--- the lemma forbids the latter. It is about the \emph{ratio} of
composition to retrieval, and about whether the derivation space is
exponentially larger than the primitive space. Under this definition:

\begin{itemize}
\tightlist
\item
  A calculator with 10 digit-arithmetic primitives that composes them
  into \(n\)-digit multiplications has \(|\Pi| = 10\) and reachable
  derivations \(\Theta(10^n)\). Ratio: exponential. \textbf{Thinker.}
\item
  A frontier LLM that has ``memorised'' (Q, A) pairs with high fidelity,
  composing at most one or two lookup steps per query, has
  \(|\Pi| \gg 10^{10}\) but \(\ell \approx 1\). Ratio: constant.
  \textbf{Pattern matcher}, regardless of parameter count.
\item
  Coq / Lean / Isabelle: \(|\Pi| \sim 10^3\) inference rules, reachable
  proofs \(\Theta(k^L)\) for branching factor \(k\) and depth \(L\).
  \textbf{Thinker.}
\item
  AlphaProof: Lean composition engine with a neural policy as the
  relevance heuristic. \textbf{Thinker} (the neural component is the
  heuristic, not the answer).
\end{itemize}

\textbf{Corollary (The Four-Part Structure of a Thinker).} Any system
with \(C(\mathcal{S}) \gg 1\) that scales as we require must have four
separable subsystems:

\begin{enumerate}
\def\labelenumi{\arabic{enumi}.}
\tightlist
\item
  \textbf{Primitives} \(\Pi\) --- small, composable, closed under the
  composition operation.
\item
  \textbf{Composition engine} --- the mechanism that assembles
  derivations. Tree search, graph rewriting, proof search, program
  synthesis. \emph{This is where the thinking happens.}
\item
  \textbf{Relevance heuristic} --- a policy that ranks which primitive
  to try next. Learned or hand-tuned. Makes composition tractable in an
  exponential space.
\item
  \textbf{Verifier} --- a decidable check that a derivation is valid.
  Non-negotiable; distinguishes construction from generation.
\end{enumerate}

The heuristic (component 3) is where neural/statistical/physics-based
systems belong. The engine (component 2) is where thinking lives.
Missing (2) reduces the system to a generator; missing (4) reduces it to
a hallucinator; missing (3) reduces it to brute-force search.

\textbf{This corrects a mistake in our own prior lineage:} every earlier
paper (CMA through RIFA) attempted to make the substrate itself the
thinker. The lemma and this corollary say that is a category error. The
substrate is the \emph{heuristic}, in the sense that AlphaProof's neural
network is a heuristic for Lean. The thinker is the composition engine
sitting on top of it. In the tri-layer architecture of §5, this is Layer
2.

The formal ratio \(C(\mathcal{S})/R(\mathcal{S})\) is the \emph{thinker
index} of a cognitive system. It is a measurable quantity, not a
philosophical claim.

\begin{center}\rule{0.5\linewidth}{0.5pt}\end{center}

\subsubsection{4. The Constructive Structural Cognition
Framework}\label{the-constructive-structural-cognition-framework}

We propose that Synthetic Intelligence is defined by three axioms
distinguishing it from statistical AI.

\paragraph{4.1 Axiom I. Categorical-Topological
State}\label{axiom-i.-categorical-topological-state}

A cognitive state is \emph{not} a vector of real numbers. A cognitive
state is a \textbf{typed graph of active relationships}: a finite
collection of nodes (concepts) connected by typed morphisms (relations),
evolving under structure-preserving transformations (functors).

Formally, let \(\mathcal{G}\) be a small category whose objects are
concept nodes and whose morphisms are typed relations. A cognitive state
is a subobject of \(\mathcal{G}\). Reasoning is the composition of
morphisms and the introduction of new subobjects via natural
transformations.

This is a strict departure from metric-space cognition. In a metric
space, similarity is a distance. In a categorical topology, similarity
is \emph{structural compatibility}: two subgraphs are similar if they
are related by a functor. The system does not measure how close two
thoughts are; it evaluates whether one can be transformed into the other
by a valid rule.

We do not require full Grothendieck-topos machinery, which is
computationally intractable at scale. In practice we use \textbf{typed
graph rewriting} as a computable specialization of the categorical
structure: nodes are concepts, edges are typed relations, and rewrite
rules are the local presentation of categorical morphisms. This
preserves the axiom's spirit while remaining implementable.

\paragraph{4.2 Axiom II. Non-Equilibrium Phase-Space
Expansion}\label{axiom-ii.-non-equilibrium-phase-space-expansion}

Classical dynamical systems have a fixed phase space determined at
construction time. A cognitive system does not. When the system
encounters a problem that its current structure cannot resolve, the
reasoning process \emph{physically spawns new degrees of freedom} ---
new nodes, new relations, new subgraphs --- to accommodate the problem.

Formally, if \(G_t\) is the graph at reasoning step \(t\), the
transition \(G_t \to G_{t+1}\) is not confined to
\(|G_{t+1}| \leq |G_t|\). New vertices and edges may be introduced by
construction rules. The Liouville conservation law that governs fixed
phase spaces does not apply; instead we introduce the
\textbf{tension-bounded expansion principle}:

\[
\Delta|G_t| \leq \alpha \cdot T(t),
\]

where \(T(t)\) is the RIFA-style tension (contradiction magnitude
weighted by log-coherence) and \(\alpha\) is a growth-rate constant.
Tension governs how much the system is \emph{permitted} to expand its
state to resolve pressure. Under stable operation (\(T\) small) the
graph stabilizes; under contradiction (\(T\) large) it grows to
accommodate.

This provides an analog of the conservation law --- not phase-space
\emph{volume} (which now grows), but the \textbf{cost of growth}, which
is charged against a tension budget. The system cannot expand
infinitely; the expansion is thermodynamically constrained.

\paragraph{4.3 Axiom III. Active Construction Over Statistical
Interpolation}\label{axiom-iii.-active-construction-over-statistical-interpolation}

A cognitive engine does not select the most likely next token from a
learned distribution. It \textbf{constructs}, per query, a deterministic
mechanism that computes the answer.

Given a query \(q\), the reasoning process assembles a subgraph
\(G^*(q) \subseteq G\) whose structure \emph{is} the answer's
derivation. The final semantic output is read from a specific node (or
invariant) of \(G^*(q)\) via the translation table.

This is the same computational primitive that underlies proof search
(Coq, Lean, Isabelle) and constraint solving (Z3, MiniSat). The
difference in CSC is that the \emph{heuristic} guiding which rewrite
rules to try --- the equivalent of a proof-tactic policy --- is provided
by the physics-based substrate rather than by a hand-crafted search
strategy or a monolithic neural policy. The substrate contributes
coherence-driven retrieval of relevant nodes and rules; the constructive
layer performs the actual derivation; the translation table names the
final result.

\paragraph{4.4 The Nine Reasoning
Primitives}\label{the-nine-reasoning-primitives}

The composition space of a constructive engine is determined not by the
\emph{number} of stored rules but by the \emph{closure} of the primitive
rule types under composition. We argue that a small finite set --- nine
primitives --- is sufficient to close the space, on the same principle
by which a finite set of grammar rules generates unbounded language
(Chomsky, 1957).

The nine primitives, and their operational form on a typed graph \(G\):

\begin{longtable}[]{@{}
  >{\raggedright\arraybackslash}p{(\linewidth - 6\tabcolsep) * \real{0.2500}}
  >{\raggedright\arraybackslash}p{(\linewidth - 6\tabcolsep) * \real{0.2500}}
  >{\raggedright\arraybackslash}p{(\linewidth - 6\tabcolsep) * \real{0.2500}}
  >{\raggedright\arraybackslash}p{(\linewidth - 6\tabcolsep) * \real{0.2500}}@{}}
\toprule\noalign{}
\begin{minipage}[b]{\linewidth}\raggedright
\#
\end{minipage} & \begin{minipage}[b]{\linewidth}\raggedright
Primitive
\end{minipage} & \begin{minipage}[b]{\linewidth}\raggedright
Pattern
\end{minipage} & \begin{minipage}[b]{\linewidth}\raggedright
Conclusion
\end{minipage} \\
\midrule\noalign{}
\endhead
\bottomrule\noalign{}
\endlastfoot
1 & \textbf{Deduction} (modus ponens) & \(P \wedge (P \Rightarrow Q)\) &
\(Q\) \\
2 & \textbf{Transitive composition} &
\((A \xrightarrow{r_1} B) \wedge (B \xrightarrow{r_2} C)\) &
\((A \xrightarrow{r_3} C)\) \\
3 & \textbf{Inheritance} (class-property flow) &
\((X \xrightarrow{\text{is\_a}} C) \wedge (C \xrightarrow{r} Y)\) &
\((X \xrightarrow{r} Y)\) \\
4 & \textbf{Abduction} & goal \(G\), rule \((P \Rightarrow G)\) &
hypothesise \(P\) \\
5 & \textbf{Negation} (contrapositive) &
\((P \Rightarrow Q) \wedge (\neg Q)\) & \((\neg P)\) \\
6 & \textbf{Case analysis} (disjunction elim.) &
\((X \text{ OR } \{d_i\}) \wedge \forall i\,(d_i \Rightarrow R)\) &
\((X \Rightarrow R)\) \\
7 & \textbf{Induction} (schematic gen.) &
\(\{X_i \text{ is\_a } C, X_i \text{ has } P\}\) for enough \(i\) &
\((C \text{ has } P)\) \\
8 & \textbf{Analogy} (structure map) & \((X \sim X')\) and
\((X \xrightarrow{r} Y)\) & \((X' \xrightarrow{r} Y)\) \\
9 & \textbf{Recursion} & any rule may fire on its own output & fixpoint
iteration \\
\end{longtable}

\textbf{Claim (Composition Closure).} Under recursive composition
(primitive 9) and the meta-rule discovery mechanism of §5.3.3, the
reachable derivation set from any consistent premise base is
Turing-complete when the ontology permits variable binding and
arithmetic. The proof is standard: the nine primitives include modus
ponens, universal instantiation (inheritance), and case analysis, which
suffice to encode any first-order proof; recursion + rule induction
allow the effective procedure to grow its rule set from experience. This
is the closure argument that makes ``hardcoded rules'' a \emph{seed}
rather than a ceiling.

\textbf{Empirical validation.} Each primitive is implemented as a
separate rule-type in \texttt{src/primitives.hpp} of the accompanying
implementation. On a small hand-crafted ontology (dogs, students,
planets, switches, chains of length 15) each primitive derives the
expected conclusion:

\begin{itemize}
\tightlist
\item
  Transitive: \texttt{dog\ is\_a\ animal} via 3 edges --- YES
\item
  Deduction: \texttt{raining\ implies\ risk} --- YES
\item
  Inheritance: \texttt{socrates\ has\ mortality},
  \texttt{socrates\ has\ language} --- YES
\item
  Negation: \texttt{socrates\ NOT\ has\_gills} via contrapositive of
  \((\text{has\_gills} \Rightarrow \text{aquatic})\) --- YES
\item
  Case analysis: \texttt{switch\_state\ ⇒\ actionable} via both
  disjuncts --- YES
\item
  Induction: \texttt{student\ has\ textbook} from 4 witnesses --- YES
\item
  Analogy: cross-domain edge transfer under a \texttt{is\_a} mapping ---
  YES
\item
  Recursion: \texttt{a\ is\_a\ p} across a 15-hop chain via fixpoint
  iteration --- YES (105 edges derived over 4 steps under tension budget
  5000)
\item
  Rule induction: macro-rule discovered from 24 forward-chaining
  derivations --- YES
\end{itemize}

Every derivation is fully traceable via the parent-edge chain (Layer 2
§5.3.4), and every conclusion was \textbf{never asserted directly}. This
is the composition-vs-retrieval distinction of §3.4 exhibited on an
actual buildable engine.

\begin{center}\rule{0.5\linewidth}{0.5pt}\end{center}

\subsubsection{5. The Tri-Layer Unified
Architecture}\label{the-tri-layer-unified-architecture}

We now specify the concrete synthesis of the three axioms with our prior
lineage.

\paragraph{5.1 Overview}\label{overview}

\begin{verbatim}
                   ┌──────────────────────────────────────────┐
INPUT (bytes) ─►   │  LAYER 1: Physics Substrate              │
                   │  ─────────────────────────               │
                   │                                           │
                   │  RFE     (kernel-field encoder)          │
                   │  DCHL Φ  (coherence potential)           │
                   │  Hamiltonian evolution                    │
                   │  RIFA IC + SIL + Crown Lattice           │
                   │                                           │
                   │  Trained via ZeroGrad Hebbian.           │
                   │  Output: (a) coherent state c*           │
                   │          (b) tension T(t)                │
                   │          (c) retrieved nodes & rules     │
                   └──────────────────┬───────────────────────┘
                                      │
                   ┌──────────────────▼───────────────────────┐
                   │  LAYER 2: Constructive Substrate         │
                   │  ─────────────────────────               │
                   │                                           │
                   │  Typed graph G of active concepts        │
                   │  Rewrite rules → morphisms G → G'        │
                   │  Growth budget bounded by T(t)           │
                   │                                           │
                   │  Reasoning = fixpoint of rule-application │
                   │  Terminates on:                          │
                   │   – answer node reached                  │
                   │   – no rule applies                      │
                   │   – tension budget exhausted             │
                   │                                           │
                   │  Realises Axioms I, II, III.             │
                   │  Output: final graph G*                  │
                   └──────────────────┬───────────────────────┘
                                      │
                   ┌──────────────────▼───────────────────────┐
                   │  LAYER 3: MDL Translation                │
                   │  ─────────────────────────               │
                   │                                           │
                   │  Table T of ≤ N rows:                    │
                   │    (graph_invariant_i, string_i)         │
                   │                                           │
                   │  Lookup: match(G*) → string              │
                   │  If no match above threshold →           │
                   │    "unknown", supported by IC's         │
                   │    reported low coherence                │
                   │                                           │
                   │  Cost: exactly log₂N bits (Lemma).       │
                   └──────────────────┬───────────────────────┘
                                      │
                                      ▼
                              SEMANTIC ANSWER
\end{verbatim}

\paragraph{5.2 Layer 1: Physics Substrate --- the Relevance
Heuristic}\label{layer-1-physics-substrate-the-relevance-heuristic}

In the four-part structure of §3.4, Layer 1 is the \textbf{relevance
heuristic}, not the thinker. Its job is the same as AlphaProof's neural
policy: given the current query and derivation state, propose
\emph{which} primitives (nodes, rules) Layer 2 should try next, in what
order. It does not answer; it ranks.

Its four responsibilities:

\begin{itemize}
\tightlist
\item
  \textbf{Semantic normalization.} Paraphrases of a query evolve to
  nearby states under the Hamiltonian, so downstream retrieval and
  rewriting see a canonical form. This is work no lookup table can do.
\item
  \textbf{Contradiction metabolism.} RIFA's Integrity Core {[}Mahir
  2026g{]} converts contradiction pressure into a scalar tension
  \(T(t)\) that regulates Layer 2's growth budget.
\item
  \textbf{Identity persistence.} RIFA's Silence-Integration Layer
  sustains the operator across long input gaps without external
  prompting.
\item
  \textbf{Rule-relevance ranking.} Given the current \((q, G_t)\),
  output a probability distribution over Layer 2's rule set --- a
  \emph{tactic policy} in the theorem-proving sense. This is where the
  substrate replaces a brute-force forward-chaining search.
\end{itemize}

The substrate is trained by the \textbf{ZeroGrad} Hebbian paradigm
{[}Mahir 2026d{]}: no backprop, no loss functions, no optimizer state,
no gradient descent on parameters. All updates are local and
physics-derived. Weight compression uses \textbf{Arch-Native
Quantization} at 4-bit precision {[}Mahir 2026f{]} for GPU-scale
substrates. \textbf{Crucially, the substrate is trained on derivation
traces from Layer 2}, not on raw text: it learns to predict which rule
was actually applied at each step of past successful derivations. This
is precisely the training signal used by AlphaZero's policy network ---
but implemented via ZeroGrad Hebbian rather than backprop.

\textbf{Anti-memorization guarantee.} The substrate never sees a
\texttt{(query,\ answer)} pair. Its training signal is
\texttt{(graph\ state,\ rule\ applied)} triples, drawn from the
derivations Layer 2 has itself generated. There is therefore no channel
through which the substrate could memorize the mapping from queries to
answers --- the mapping is not in the training data. It could, at worst,
memorize \texttt{(state,\ rule)} correlations, which is exactly what a
heuristic policy is \emph{supposed} to do (see §5.2's list of
responsibilities). The engine's answers come from Layer 2's construction
plus Layer 3's translation table; neither of these paths involves the
substrate as anything other than a proposer of what to try next. This is
what closes the ``hidden memorization'' trap that
pattern-matcher-in-thinker's-clothing systems fall into.

\paragraph{5.3 Layer 2: Constructive Substrate --- the
Thinker}\label{layer-2-constructive-substrate-the-thinker}

Layer 2 is the composition engine of §3.4. This is where thinking lives.
It is a \textbf{typed graph rewriter} operating on a graph
\(G_t = (V_t, E_t, \tau)\) where \(\tau: E_t \to \mathcal{R}\) assigns
each edge a relation type from a finite ontology \(\mathcal{R}\). Its
four subcomponents mirror the corollary:

\textbf{5.3.1 Rule Family.} Rewrite rules are of the form

\[
\text{if } (P \subseteq G_t) \text{ then } G_{t+1} := G_t \cup Q,
\]

where \(P\) is a pattern subgraph and \(Q\) is a construction subgraph.
New nodes in \(Q \setminus P\) realize Axiom II (phase-space expansion).
The initial rule set is the \textbf{nine reasoning primitives of §4.4}
(deduction, transitive composition, inheritance, abduction, negation,
case analysis, induction, analogy, recursion). Each is implemented as a
distinct rule type with its own pattern matcher; the accompanying
implementation \texttt{src/primitives.hpp} fires all nine on the
empirical suite of §4.4.

By the composition-closure claim of §4.4, this seed set is sufficient to
reach a Turing-complete derivation space under recursion + rule
induction. The set is fixed at design time; the \emph{macros} discovered
from usage (§5.3.3) multiply the reachable derivation space
monotonically over the system's lifetime, without new hand-authored
rules.

\textbf{5.3.2 Composition Engine.} Both forward and backward chaining.
Forward derives everything currently reachable; backward pursues a
specific query goal by working from the goal to the premises. The engine
is a bounded-depth iterative deepening: forward for exploration,
backward for goal-directed proof. Termination conditions: (a) answer
node reached, (b) no rule applies (fixpoint), (c) tension-bounded budget
\(\alpha \cdot T(t)\) exhausted.

\textbf{5.3.3 Rule Induction (Meta-Rule Discovery).} After every
successful derivation, sub-chains of length \(\geq 2\) are recorded with
their frequency. When a sub-chain crosses a MDL threshold (its
compressed representation is shorter than the sum of its parts), it is
registered as a new \textbf{macro-rule}. Subsequent derivations that
would have replayed the sub-chain now apply the macro-rule instead. This
is how the system grows its own primitive set from experience --- the
primitive count \(P\) grows polynomially, but each new macro-rule
multiplies the reachable derivation space. The result is a monotone
increase in the thinker index \(C/R\) over the system's lifetime, purely
from usage, without any additional training data.

\textbf{5.3.4 Verifier.} Every derived edge carries the rule that
produced it and the parent edges. A verifier walks any candidate answer
edge back to premises and re-checks each application. This is the
mechanism that prevents Layer 2 from being a generator: no output is
emitted unless its derivation type-checks against the current rule set.
Failed verification returns ``unknown'' via Layer 3, not a hallucinated
answer.

This layer is what none of our prior work possessed. It is what turns
coherent-state normalisation into constructive derivation, and ---
through §5.3.3 --- enables the system to \emph{increase its thinker
index over time} rather than being frozen at design.

\paragraph{5.4 Layer 3: Minimum-Description-Length
Translation}\label{layer-3-minimum-description-length-translation}

The translation table pays the conservation-law tax explicitly. Each row
is a pair \((I_i, s_i)\) where \(I_i\) is a graph invariant (e.g., the
label of a designated answer node, or a signature over \(G^*\)) and
\(s_i\) is its human-readable string. Lookup is exact (invariants are
discrete). The total cost is \(\Theta(N)\) table rows for \(N\) answers,
meeting the \(\log_2 N\) lower bound within a constant factor.

Because the translation is explicit, it is \textbf{interpretable} in a
way no distributed LLM head is: every row can be inspected, every answer
traced to a specific graph structure. Because the invariants are
discrete, translation \textbf{does not hallucinate}: an answer is either
a table hit or ``unknown'' (reported with IC-supported low coherence).

\paragraph{\texorpdfstring{5.4b Physics-Driven Language Generation via
Resonance Signatures ---
\textbf{RETIRED}}{5.4b Physics-Driven Language Generation via Resonance Signatures --- RETIRED}}\label{b-physics-driven-language-generation-via-resonance-signatures-retired}

\begin{quote}
\textbf{STATUS: RETIRED. Superseded by §5.4d (Constructive Language
Locking).}

This section is preserved because it states a hypothesis clearly and
that hypothesis was \textbf{wrong in kind, not in degree} --- which is a
result worth keeping. Do not cite it as describing the current system.

The proposal below asks a resonance field to \emph{discover} English
word order. §5.4d argues that this is a category error: grammar is a
transmission protocol two speakers agreed on, not a physical phenomenon
awaiting measurement. There is no gradient to follow because there is no
fact of the matter. The concept-level ordering described in Stage 3
partially worked and survives in CPME/CPSB; the \emph{grammatical}
claims of Stage 2 did not, and the byte-level emission babbled
(\texttt{mrs.md} §4b).

\textbf{What replaced it:} surface realisation moved out of the physics
entirely, into isolated deterministic language clusters (§5.4d,
\texttt{cll.md}). What was retained: phase-collapse ordering across
\emph{concepts}, which is a physical claim and holds.

The comparison table at the end of this section should be read as
\emph{what was hoped for in 2026}, not as a measured result. Two of its
rows (``Grammar: fixed macro-node ω per grammatical role'',
``Sequencing: phase-locked slot sampling'') describe a mechanism that
was implemented and did not work.
\end{quote}

The MDL translation table of §5.4 answers a single query with a single
string via a static lookup on a graph invariant. Natural-language
\emph{generation} --- producing ordered sentences --- requires an
extension: the wave dynamics of Layer 1 already encode meaning, but the
\emph{sequence} of words must emerge from a mechanism that does not
sample from a next-token distribution. We proposed a three-stage
physics-only pipeline that produces grammatically ordered output without
any statistical prediction.

\textbf{Stage 1 --- Resonance Signature (semantics).} A query drives the
cognitive substrate; after a settling window the standing-wave field
\(\{x_i, v_i, \theta_k\}\) across cavities and macro-nodes IS the
meaning. This is consistent with the fixed-point account of phenomenal
consciousness in BSMA {[}Mahir 2026m, \texttt{bmsa\_new.md}{]}: the wave
field is not a proxy for a thought --- it is the thought. No decoding is
required at this stage --- the meaning lives in the physical state, and
the primitives it references live in the graph.

\textbf{Stage 2 --- Phase-Locking with a Syntax Macro-Layer.} In
addition to the semantic macro-nodes of HNCA {[}Mahir 2026n,
\texttt{hnca.md}{]} we introduce a second macro-layer whose regions have
\emph{grammatical roles}: one macro-node per slot (Subject, Verb,
Object, and optional Adverb / Modifier) with a fixed natural frequency
\(\omega_S, \omega_V, \omega_O\). The semantic macro-nodes couple into
the syntax layer via the same Kuramoto rule already used for HNCA:

\[
\frac{d\theta_i}{dt} = \omega_i + \frac{K}{M}\sum_j \sin(\theta_j - \theta_i)\,A_j
\]

When a semantic cluster's amplitude peaks near the Subject slot's phase,
that cluster IS the subject of the emitted sentence. This is the same
mechanism a human uses when a ``verb slot'' in Broca's area recruits
whichever motor-action concept currently has the highest neural gain ---
grammatical role is a \emph{physical slot}, not a statistical
association.

\textbf{Stage 3 --- Temporal Readout Head with Tunnel Decoding.} A
readout head samples the cavities at each slot phase in canonical order:
\(\theta_S = 0\), \(\theta_V = 2\pi/3\), \(\theta_O = 4\pi/3\). For each
slot the head reports the cavity whose displacement is peak \emph{at
that phase}. Only then are the raw label strings concatenated. The
result is not a bag of concepts --- it is an ordered sentence whose
sequence is determined by physics, not by probability. Concretely, for a
probe on \texttt{metformin}:

\begin{itemize}
\tightlist
\item
  The semantic wave excites \texttt{metformin.α}, which phase-locks into
  the Subject slot; the tunnel
  \texttt{metformin\ -\/-inhibits-\/-\textgreater{}\ mtor} transmits
  into the Verb slot at phase \(2\pi/3\); the target cavity
  \texttt{mtor.β} phase-locks into the Object slot at \(4\pi/3\).
\item
  The readout head samples these three slots in the fixed canonical
  order.
\item
  Layer-3 table maps each cavity id to its label string.
\item
  Output: \texttt{metformin\ inhibits\ mtor} --- one sentence,
  deterministic, physics-derived, no next-token prediction anywhere in
  the pipeline.
\end{itemize}

\textbf{Storage-floor honesty.} The three raw label strings
(\texttt{metformin}, \texttt{inhibits}, \texttt{mtor}) ARE stored ---
they must be, by the Translation Capacity Lemma of §3.1. The
\emph{sequencing} is NOT stored. No pair
\texttt{(query,\ "metformin\ inhibits\ mtor")} appears anywhere in the
training data. The order emerges from the physical geometry of the
syntax macro-layer coupled into the semantic substrate. This is the same
storage floor a human brain requires: a child who has never heard
\texttt{metformin} cannot dream about it. Facts must be stored;
sentences need not be.

\textbf{Comparison.}

\begin{longtable}[]{@{}
  >{\raggedright\arraybackslash}p{(\linewidth - 4\tabcolsep) * \real{0.3333}}
  >{\raggedright\arraybackslash}p{(\linewidth - 4\tabcolsep) * \real{0.3333}}
  >{\raggedright\arraybackslash}p{(\linewidth - 4\tabcolsep) * \real{0.3333}}@{}}
\toprule\noalign{}
\begin{minipage}[b]{\linewidth}\raggedright
Mechanism
\end{minipage} & \begin{minipage}[b]{\linewidth}\raggedright
LLM
\end{minipage} & \begin{minipage}[b]{\linewidth}\raggedright
This engine
\end{minipage} \\
\midrule\noalign{}
\endhead
\bottomrule\noalign{}
\endlastfoot
Semantics & Distributed in \(\sim 10^{12}\) weights & Standing-wave
field on the substrate \\
Sequencing & \(\arg\max\) over next-token distribution & Phase-locked
slot sampling \\
Grammar & Learned from token co-occurrence & Fixed macro-node \(\omega\)
per grammatical role \\
Sampling & Temperature \(\times\) softmax & None --- deterministic slot
readout \\
Novelty & Statistical interpolation & Physics-computed constructions
(fission, bridge insertion, resonance) \\
Storage & \((Q,A)\) pairs baked into weights & Only primitive fact
strings; \(O(N)\) rows \\
Failure mode & Hallucination & Explicit ``unknown'' via slot-vacancy \\
\end{longtable}

This mechanism is the natural extension of Layer 3 --- the same MDL
table supplies the label strings; the syntax macro-layer supplies the
ordering that a single-slot lookup cannot. Neither statistical sampling
nor an autoregressive language model appears anywhere in the pipeline.
The extension is a small addition: one macro-layer with three-to-five
fixed frequencies, one readout head, no new weights, no new training.

\paragraph{5.4c The Unified Working Substrate --- What Is Built,
Measured, and
Load-Bearing}\label{c-the-unified-working-substrate-what-is-built-measured-and-load-bearing}

The papers in this lineage each propose one mechanism. Read separately
they are a pile of proposals; read as a stack they are one substrate in
which each mechanism gates a different physical quantity. This section
arranges the mechanisms that are \textbf{implemented and measured} ---
not proposed --- in the order the field encounters them. Everything
below is in \texttt{src/physics.hpp} unless noted.

\textbf{Input path.}

\begin{longtable}[]{@{}
  >{\raggedright\arraybackslash}p{(\linewidth - 6\tabcolsep) * \real{0.2500}}
  >{\raggedright\arraybackslash}p{(\linewidth - 6\tabcolsep) * \real{0.2500}}
  >{\raggedright\arraybackslash}p{(\linewidth - 6\tabcolsep) * \real{0.2500}}
  >{\raggedright\arraybackslash}p{(\linewidth - 6\tabcolsep) * \real{0.2500}}@{}}
\toprule\noalign{}
\begin{minipage}[b]{\linewidth}\raggedright
Mechanism
\end{minipage} & \begin{minipage}[b]{\linewidth}\raggedright
Paper
\end{minipage} & \begin{minipage}[b]{\linewidth}\raggedright
Gates
\end{minipage} & \begin{minipage}[b]{\linewidth}\raggedright
Status
\end{minipage} \\
\midrule\noalign{}
\endhead
\bottomrule\noalign{}
\endlastfoot
Byte-lane sensory layer & CMD & 256 lanes; concepts reached only via
learned bindings & built \\
Bigram lanes & MRS §4b.1 & ordered pairs; recognition survives
misspelling & \textbf{7/7 measured} \\
Chirality & CPME §1.1 & arrival time → phase gradient & built \\
\end{longtable}

\textbf{Field dynamics.} Each row multiplies a different term; none is a
special case of another.

\begin{longtable}[]{@{}
  >{\raggedright\arraybackslash}p{(\linewidth - 6\tabcolsep) * \real{0.2500}}
  >{\raggedright\arraybackslash}p{(\linewidth - 6\tabcolsep) * \real{0.2500}}
  >{\raggedright\arraybackslash}p{(\linewidth - 6\tabcolsep) * \real{0.2500}}
  >{\raggedright\arraybackslash}p{(\linewidth - 6\tabcolsep) * \real{0.2500}}@{}}
\toprule\noalign{}
\begin{minipage}[b]{\linewidth}\raggedright
Mechanism
\end{minipage} & \begin{minipage}[b]{\linewidth}\raggedright
Paper
\end{minipage} & \begin{minipage}[b]{\linewidth}\raggedright
Gates
\end{minipage} & \begin{minipage}[b]{\linewidth}\raggedright
Physical quantity
\end{minipage} \\
\midrule\noalign{}
\endhead
\bottomrule\noalign{}
\endlastfoot
MIMS materials & \texttt{mims.md} & tunnel stiffness / damping &
\textbf{now derived from fan-out, not hand-set} \\
Curie habituation & CMD & tunnel coupling & per-cavity heat \(T_i\) \\
VGHT sub-bands & \texttt{vght.md} & domain isolation & macro-node
\(\omega\) \\
VGHT voltage gates & \texttt{vght.md} & cross-domain bridges &
\(\sigma(\kappa(|x_a||x_b| - \theta))\) \\
VGHT viscosity & \texttt{vght.md} & avalanche suppression &
\(E_{\text{total}}\) \\
CPSB gauge \(\Psi\) & \texttt{cpsb.md} & directed phase order &
per-tunnel offset \\
CSAT chaos & \texttt{CSAT.md} & metastability & Lorenz
\(\hat x \cdot \tanh(T_i/T_c)\) \\
Energy homeostasis & this §5.4c & inter-turn forgetting & global gain
\(\sqrt{B/E}\) \\
\end{longtable}

\textbf{Structural mechanisms.}

\begin{longtable}[]{@{}
  >{\raggedright\arraybackslash}p{(\linewidth - 4\tabcolsep) * \real{0.3333}}
  >{\raggedright\arraybackslash}p{(\linewidth - 4\tabcolsep) * \real{0.3333}}
  >{\raggedright\arraybackslash}p{(\linewidth - 4\tabcolsep) * \real{0.3333}}@{}}
\toprule\noalign{}
\begin{minipage}[b]{\linewidth}\raggedright
Mechanism
\end{minipage} & \begin{minipage}[b]{\linewidth}\raggedright
Paper
\end{minipage} & \begin{minipage}[b]{\linewidth}\raggedright
Effect
\end{minipage} \\
\midrule\noalign{}
\endhead
\bottomrule\noalign{}
\endlastfoot
RB+G fission & \texttt{rbg.md} & contradiction → new vertices (Axiom
II) \\
L6-A bridging & \texttt{rbg.md} §5 & shared-target intersection → rubber
tunnel \\
ZeroGrad Hebbian & \texttt{zerograd.md} & local plasticity, no
backprop \\
CPME-II phase Hebbian & \texttt{cpme\_ii.md} & consolidates only
phase-consistent co-activity \\
\end{longtable}

\textbf{Two corrections this arrangement forced.} Assembling the stack
revealed defects invisible when each paper was read alone:

\begin{enumerate}
\def\labelenumi{\arabic{enumi}.}
\item
  \emph{MIMS constants were guesses.} Stiffness is now \textbf{derived}
  from a relation's fan-out:
  \(k_r = k_{hi}/(1 + \log_2(1 + \text{fanout}_r))\). A relation whose
  subject almost determines its object is a hard constraint and
  transmits cleanly; a promiscuous one is diffuse and should not. Twenty
  hand-set numbers became one formula, and the physics measurably
  improved.
\item
  \emph{VGHT sub-bands were assigned per connected component, which
  silently failed.} Union-find over \textbf{all} tunnels merged 69\% of
  the graph (1877 of 2712 cavities) into one component, which then
  received a single frequency band --- so domain isolation was absent
  exactly where it was needed, and biology, taxonomy and computing
  resonated in the same band. The cause was generic \texttt{is\_a} /
  \texttt{has} edges bridging every domain through shared abstract hubs.
  The fix follows from correction (1): those relations now carry
  \emph{low} stiffness by construction, so clustering unions only over
  tunnels above median stiffness. \textbf{Two masses joined by a very
  weak spring are not one rigid body.} The giant component fell to 621
  cavities and cross-domain contamination disappeared:

\begin{verbatim}
before   mammal -> biology open source license kernel posix compliance
after    mammal -> warm blood hair spine
\end{verbatim}
\end{enumerate}

\textbf{Verification path (Phase 2).} The substrate now forms hypotheses
from the graph, emits real source expressing them, and submits them to
an external oracle:

\begin{longtable}[]{@{}
  >{\raggedright\arraybackslash}p{(\linewidth - 2\tabcolsep) * \real{0.5000}}
  >{\raggedright\arraybackslash}p{(\linewidth - 2\tabcolsep) * \real{0.5000}}@{}}
\toprule\noalign{}
\begin{minipage}[b]{\linewidth}\raggedright
Fact
\end{minipage} & \begin{minipage}[b]{\linewidth}\raggedright
Provenance
\end{minipage} \\
\midrule\noalign{}
\endhead
\bottomrule\noalign{}
\endlastfoot
\texttt{human\ is\_a\ great\_ape} & derived --- forward chaining,
premises cited \\
\texttt{bind\_byte\ has\ arity\_2} & \textbf{established by experiment}
--- contradicted extraction; compiler settled it \\
\texttt{drive\ has\ sig\_int\_float\_float\_int} & \textbf{established
by experiment} --- absent from the graph entirely \\
\end{longtable}

This adds a third provenance class to the substrate. Neither asserted
nor derived: \emph{established by experiment}. Cost on a repeated task
falls from 3 experiments to 1, because the hypothesis it starts from is
the one it earned.

\paragraph{5.4d Language Governance --- Isolated Syntax
Clusters}\label{d-language-governance-isolated-syntax-clusters}

\emph{Developed at length in \texttt{research-paper/cll.md}
(Constructive Language Locking). This section states the principle and
the measured result; the companion paper gives the formal definition,
cost accounting, comparison table and failure modes.}

§5.4b proposed deriving word order from field physics, and §8.5 asked
for ``grammatical constraint expressed as a field property.''
\textbf{That was a category error, and this section corrects it.}

\textbf{Cognition and language are different kinds of thing.}

\begin{longtable}[]{@{}
  >{\raggedright\arraybackslash}p{(\linewidth - 4\tabcolsep) * \real{0.3333}}
  >{\raggedright\arraybackslash}p{(\linewidth - 4\tabcolsep) * \real{0.3333}}
  >{\raggedright\arraybackslash}p{(\linewidth - 4\tabcolsep) * \real{0.3333}}@{}}
\toprule\noalign{}
\begin{minipage}[b]{\linewidth}\raggedright
\end{minipage} & \begin{minipage}[b]{\linewidth}\raggedright
Cognition
\end{minipage} & \begin{minipage}[b]{\linewidth}\raggedright
Language
\end{minipage} \\
\midrule\noalign{}
\endhead
\bottomrule\noalign{}
\endlastfoot
Purpose & discover & transmit \\
Needs & novelty, experiment, resonance, oracle validation & a fixed,
agreed protocol \\
Bound by rules? & never --- that would cripple it & entirely --- that is
what it \emph{is} \\
\end{longtable}

English SVO is not a physical law awaiting discovery by resonance. It is
a human communication convention. Bengali SOV encodes identical
cognitive content under a different convention. Attempting to make the
substrate \emph{discover} grammar confuses a protocol with a phenomenon.

\textbf{Architecture.} The substrate splits, and the core is untouched:

\begin{verbatim}
[ Core SI Physics Substrate ]  ->  raw resonant concept triples
                |                  (subject, relation, object)
                v
[ Isolated Language Cluster ]  ->  deterministic surface realization
     |- English  (SVO grid)
     |- Bengali  (SOV grid)
     |- Banglish (SOV grid, Latin script)
                |
                v
        human-readable stream
\end{verbatim}

Not one grammar rule enters \texttt{physics.hpp}. Fission, L6-A
bridging, thermodynamic discovery and novelty are unaffected --- the
cognitive flexibility the rest of this paper argues for is preserved
\emph{because} the protocol lives outside the physics.

\textbf{Measured} (\texttt{src/lang.hpp}), one cognitive state under
three protocols:

\begin{verbatim}
(biology, is_a, science)     en: Biology is a science.
                             bn: Biology science একটি।
                             bl: Biology science ekta.

(metformin, inhibits, mtor)  en: Metformin inhibits mtor.
                             bn: Metformin mtor বাধা দেয়।
                             bl: Metformin mtor badha dey.
\end{verbatim}

\textbf{Cost accounting.} This adds no hidden cost and hides nothing.
Layer 3 already stores label strings --- the Translation Capacity Lemma
(§3.1) forces it. A word-order constant is \textasciitilde2 bits; a
relation lexicalisation table is one row per relation per language; a
transliteration map is one row per word pair. All are Layer-3
\emph{protocol} cost, explicitly named, and none is intelligence.

\textbf{A defect testing disproved immediately.} The first
implementation treated article insertion as a global language property
and emitted \emph{``Biology is a the science.''} Whether an object takes
a determiner belongs to the \textbf{construction}, not the language:
\texttt{is\_a} supplies its own article, proper nouns take none. Article
placement is now a per-relation flag. This is the shape of the remaining
work --- real grammar exceeds word order, but every additional rule is
deterministic data, never a learned parameter.

\textbf{Two properties this buys that a language model cannot.}

\begin{enumerate}
\def\labelenumi{\arabic{enumi}.}
\tightlist
\item
  \textbf{Zero-hallucination surface form.} The transliteration relation
  is strict and static: \texttt{tumi} is \emph{always} \texttt{তুমি}. A
  model code-switching Banglish drifts between spellings because its
  output is sampled; this cannot, because nothing is sampled.
\item
  \textbf{O(1) multilingual swap.} Adding Japanese means adding a table.
  The same cognitive state emits Japanese with the substrate unchanged.
  A language model requires retraining or far greater capacity for the
  same capability.
\end{enumerate}

\textbf{Still open in this layer} --- all deterministic, none requiring
ML: morphological agreement (\emph{studies}, not \emph{study}), Bengali
verb conjugation by person and honorific (তুমি / আপনি / তুই), free word
order under case marking, and multi-triple discourse with pronoun
resolution.

\paragraph{5.4e The J-A-S Cycle --- Einstein's Epistemology as Substrate
Architecture}\label{e-the-j-a-s-cycle-einsteins-epistemology-as-substrate-architecture}

Einstein described the structure of physical theory as a cycle rather
than a ladder:

\begin{verbatim}
                  +---------------------------+
                  |        Axioms (A)         |
                  +---------------------------+
                     /                     \
                    v                       v
                +-----+                   +-----+
                |  S  |                   | S'  |   deduced theorems
                +-----+                   +-----+
                   |                         |
                   | experiment              | experiment
                   v                         v
~~~~~~~~~~~~~~~~~~~~~~~~~~~~~~~~~~~~~~~~~~~~~~~~~~~~~
                Sense Experience (E)
~~~~~~~~~~~~~~~~~~~~~~~~~~~~~~~~~~~~~~~~~~~~~~~~~~~~~
        \                                     /
         +------------- Jump (J) ------------+
\end{verbatim}

The crucial claim is that \textbf{J is not logical}. No inference rule
takes you from observations to axioms; the leap is inductive,
unjustified, and revisable. Deduction runs \emph{downward} from axioms
to theorems; experiment tests theorems back against experience;
refutation propagates \emph{upward} to revise the axioms that produced
them.

This is not an analogy for the architecture --- it is a description of
it. Every element is already implemented:

\begin{longtable}[]{@{}
  >{\raggedright\arraybackslash}p{(\linewidth - 4\tabcolsep) * \real{0.3333}}
  >{\raggedright\arraybackslash}p{(\linewidth - 4\tabcolsep) * \real{0.3333}}
  >{\raggedright\arraybackslash}p{(\linewidth - 4\tabcolsep) * \real{0.3333}}@{}}
\toprule\noalign{}
\begin{minipage}[b]{\linewidth}\raggedright
Einstein
\end{minipage} & \begin{minipage}[b]{\linewidth}\raggedright
Substrate
\end{minipage} & \begin{minipage}[b]{\linewidth}\raggedright
Implementation
\end{minipage} \\
\midrule\noalign{}
\endhead
\bottomrule\noalign{}
\endlastfoot
\textbf{E} --- sense experience & raw source, byte input &
\texttt{extract\_ast\_axioms.py}, byte-lane layer \\
\textbf{J} --- the jump & induction: universal generalisation over
observations & §5.3.3 rule induction, L6-A bridging, RB+G fission \\
\textbf{A} --- axioms & the typed concept graph &
\texttt{data/*.axioms} \\
\textbf{S, S'} --- deduced theorems & forward chaining & 100 derived
edges, premises cited \\
\textbf{experiment} & compile-verify against an external oracle &
\texttt{synth}, §5.4c \\
\textbf{refutation → revise A} & corrected axiom written back &
\texttt{data/verified.axioms} \\
\end{longtable}

\textbf{Why this matters for the SI claim.} §3.4 defines a thinker by
composition rate. The J-A-S cycle sharpens that: a system that only
deduces (S from A) is a theorem prover, and a system that only observes
(E) is a sensor. What makes something a \emph{scientist} is the closed
loop --- jumping to axioms that are not entailed, deducing their
consequences, and letting an external oracle refute them.

The substrate closes that loop. It is the strongest available answer to
the charge that this is an indexing engine: an index has no J, no
experiment, and no mechanism by which reality can overturn what it
holds.

\textbf{The essential asymmetry.} J is \emph{not} the inverse of
deduction and must never be implemented as such. An induced axiom is a
guess whose only warrant is survival under experiment. The substrate
therefore tracks three provenance classes --- \emph{asserted},
\emph{derived}, \emph{established by experiment} --- and a fourth is
added here: \emph{induced}, which carries the weakest warrant and is the
first thing retracted under refutation.

\paragraph{5.5 What Each Prior Component Contributes, and What Was
Dropped}\label{what-each-prior-component-contributes-and-what-was-dropped}

\begin{quote}
\textbf{CORRECTED.} An earlier version of this table listed DCHL, BED,
ANQ and CMC as \emph{retained components of the current architecture}.
They are not, and were not at the time the claim was written. The table
below is reconciled against the mechanism-level audit in §6.2, which was
produced by grepping the live source rather than by reading the papers.
Where the two disagreed, the grep won.

The error is worth naming because it is the characteristic failure of a
long lineage: a synthesis section written once and never re-checked
silently becomes a description of a system that no longer exists. That
is the failure §6.2 exists to prevent, and it had already happened here.
\end{quote}

The substrate this paper describes is the \textbf{oscillator substrate}
(\texttt{src/physics.hpp}). It replaced the \textbf{DCHL substrate} ---
a continuous-state Hamiltonian system with a learned potential \(\Phi\)
--- which is retired. Everything below is stated relative to that
replacement.

\begin{longtable}[]{@{}
  >{\raggedright\arraybackslash}p{(\linewidth - 4\tabcolsep) * \real{0.3333}}
  >{\raggedright\arraybackslash}p{(\linewidth - 4\tabcolsep) * \real{0.3333}}
  >{\raggedright\arraybackslash}p{(\linewidth - 4\tabcolsep) * \real{0.3333}}@{}}
\toprule\noalign{}
\begin{minipage}[b]{\linewidth}\raggedright
Component
\end{minipage} & \begin{minipage}[b]{\linewidth}\raggedright
Status in the CURRENT engine
\end{minipage} & \begin{minipage}[b]{\linewidth}\raggedright
Where, or why not
\end{minipage} \\
\midrule\noalign{}
\endhead
\bottomrule\noalign{}
\endlastfoot
CMA fixed potential & \textbf{historical only} & supplied the manifold
vocabulary of §4; no code \\
P-CMA predictive transition & \textbf{dropped} & subsumed by
graph-rewriter transitions \\
HCMA hierarchical states & \textbf{dropped with DCHL} & its contrastive
training is the negative result of §2.1 \\
CCH curl-free constraint & \textbf{not implemented} & grep finds no
\texttt{curl} in the substrate; the constraint was never enforced in
code, so it guarantees nothing here \\
DCHL learned \(\Phi\) & \textbf{RETIRED} & the substrate that collapsed
into the §1.2 traps; trained artifacts archived out of the tree \\
\textbf{ZeroGrad learning rule} & \textbf{LIVE --- foundational} &
\texttt{plasticity\_delta\ =\ pd·hebbian\_decay\ +\ plastic\_rate·v\_a·v\_b};
every Hebbian update, including CPME-II, is this rule. \emph{The paper
is dead; the rule is load-bearing} (§6.2.4) \\
\textbf{RIFA IC/SIL} & \textbf{LIVE} & \texttt{IntegrityCore},
\texttt{SilenceIntegration} --- contradiction and identity metabolism \\
RIFA Rule Synthesizer & \textbf{live but bypassed on bulk load} &
streaming-only; feeding it a sorted axiom file manufactures rules from
alphabetical order, so ingestion disables it \\
BED (Byte Energy Descent) & \textbf{DEAD} &
\texttt{DEAD/deprected/bed.md}; converged to byte-mean with no prompt
structure \\
CognitiveMemory retrieval & \textbf{dropped} & subsumed by the Layer-3
table, which names the cost honestly \\
STAN topological readout & \textbf{dropped} & bifurcation drift;
replaced by graph invariants \\
ANQ 4-bit quantisation & \textbf{DEAD as a scheme} & retired with the
GPU trainer it served \\
\textbf{ANQ LEA integrator} & \textbf{LIVE} & continuous integrator
replacing threshold-commit; used by every saturating accumulator in the
tick (§6.2.4) \\
CMC memory reduction & \textbf{not in live code} & no current trainer
needs it --- there is no gradient training left to compress \\
\textbf{LTO} & \textbf{LIVE} & runtime axiom acquisition and
persistence \\
\end{longtable}

Two entries above are the point of the whole audit: \textbf{ZeroGrad and
ANQ-LEA are running in production while their papers sit in
\texttt{DEAD/}.} A dead paper is not a retracted one. Cite those two as
\emph{sources of mechanisms}, never as \emph{descriptions of the
system}.

The architecture is minimal in the sense that removing any \textbf{LIVE}
component costs a specific capability. It is \emph{not} minimal in the
sense that every historical component found a home --- most did not, and
saying otherwise was the error corrected above.

\begin{center}\rule{0.5\linewidth}{0.5pt}\end{center}

\subsubsection{6. Relationship to Prior Work in the
Lineage}\label{relationship-to-prior-work-in-the-lineage}

\begin{quote}
\textbf{Read this section as intellectual history, not as a parts list.}
Every item below states what a paper \emph{contributed to the thinking}.
Whether the mechanism survives into running code is a separate question,
answered --- per paper, by grep --- in §6.2. Several entries here are
historically important and currently dead (CCH, CMC, ANQ, SCER, BED).
Conflating ``this shaped the design'' with ``this is in the system'' is
precisely the error corrected in §5.5.
\end{quote}

The synthesis makes explicit what each earlier paper contributed, and
what none of them could have done alone:

\begin{itemize}
\tightlist
\item
  \textbf{CMA} (2026a) introduced the manifold and Hamiltonian
  formalism.
\item
  \textbf{P-CMA} (2026b) added learned transitions and the planning
  idea.
\item
  \textbf{CCH} (2026c) established the curl-free constraint as the
  mathematical guarantor of transitive reasoning.
\item
  \textbf{HCMA} (2026h) introduced decoupled contrastive Hebbian
  learning (DCHL) and the two-level state.
\item
  \textbf{ZeroGrad} (2026d) formalized the local, backprop-free training
  paradigm.
\item
  \textbf{CMC} (2026i) reduced training memory to \(O(1)\) via cognitive
  state compression.
\item
  \textbf{ANQ} (2026f) achieved 4-bit weight compression that preserves
  DCHL physics.
\item
  \textbf{CME} (2026j) extended the substrate to multimodal inputs via
  unified state.
\item
  \textbf{LTO} (2026k) introduced online substrate adaptation without
  stopping inference.
\item
  \textbf{SCER} (2026l) proposed multi-mode routing at the reasoning
  layer.
\item
  \textbf{RIFA} (2026g) added structural pressure, contradiction
  metabolism, silence endurance, and self-assembling control.
\item
  \textbf{BED} (2026e) introduced the discrete argmin decoder for
  byte-level output.
\end{itemize}

Each was locally correct within its scope. The synthesis reveals that
none was sufficient alone to escape the three traps of §1.2, and that
the missing piece --- the constructive graph-rewriting layer with
explicit MDL translation --- is precisely what turns \emph{substrate}
into \emph{thinker}.

\textbf{Two corrections to the list above, both established by §6.2.}
CCH's curl-free constraint is described here as having ``established'' a
guarantor of transitive reasoning; it was never implemented, so it
guarantees nothing about this engine and the claim is historical only.
And the framing ``each was locally correct'' is too generous to HCMA and
BED: HCMA's contrastive objective was \emph{wrong} in the specific way
§2.1 diagnoses (it trains a classifier), and BED's decoder demonstrably
converged to the byte-mean. Those are negative results, and they are
more useful stated as such.

\begin{center}\rule{0.5\linewidth}{0.5pt}\end{center}

\paragraph{6.2 Lineage Audit --- What Still Carries Weight, and What Is
Dead}\label{lineage-audit-what-still-carries-weight-and-what-is-dead}

A lineage of thirty-plus papers accumulates claims faster than it
retires them. This section states, for every prior paper, whether it is
load-bearing in the current engine. The audit was performed by grepping
the live source for each mechanism, not by reading the papers.

\textbf{The distinction that matters most: a paper can be dead while its
mechanism survives.} Research lineages do not retire ideas cleanly; they
cannibalise them. Marking a paper dead says the \emph{proposal as
stated} is retired. It does not say the idea was worthless, and in two
cases below the idea is running in production while its paper sits in
\texttt{DEAD/}.

\subparagraph{\texorpdfstring{6.2.1 Load-bearing today
(\texttt{research-paper/})}{6.2.1 Load-bearing today (research-paper/)}}\label{load-bearing-today-research-paper}

Mechanisms verified present in \texttt{src/physics.hpp},
\texttt{src/layer2.hpp}, or the live binaries:

\begin{longtable}[]{@{}
  >{\raggedright\arraybackslash}p{(\linewidth - 4\tabcolsep) * \real{0.3333}}
  >{\raggedright\arraybackslash}p{(\linewidth - 4\tabcolsep) * \real{0.3333}}
  >{\raggedright\arraybackslash}p{(\linewidth - 4\tabcolsep) * \real{0.3333}}@{}}
\toprule\noalign{}
\begin{minipage}[b]{\linewidth}\raggedright
Paper
\end{minipage} & \begin{minipage}[b]{\linewidth}\raggedright
Mechanism
\end{minipage} & \begin{minipage}[b]{\linewidth}\raggedright
Where
\end{minipage} \\
\midrule\noalign{}
\endhead
\bottomrule\noalign{}
\endlastfoot
\texttt{prm.md} & mass-spring cavities & \texttt{struct\ Cavity} \\
\texttt{mims.md} & material impedance & \texttt{material\_for()} ---
\textbf{now derived, not hand-set} \\
\texttt{hnca.md} & macro Kuramoto & \texttt{struct\ MacroNode} \\
\texttt{ntda.md} & neuromodulators &
\texttt{struct\ Neurotransmitters} \\
\texttt{tsda.md} & salience integrator & \texttt{Cavity::salience} \\
\texttt{rifa.md} & contradiction + identity & \texttt{IntegrityCore},
\texttt{SilenceIntegration} \\
\texttt{srg.md} & dual-port routing & \texttt{x\_E/x\_I} (opt-in) \\
\texttt{rbg.md} & gated fission & \texttt{fission\_locked()} \\
\texttt{vght.md} & sub-bands, voltage gates & band ladder,
\texttt{voltage\_gated} \\
\texttt{cpme.md}, \texttt{cpme\_ii.md} & momentum emission, phase
Hebbian & \texttt{readout\_phase\_collapse()} \\
\texttt{cpsb.md} & gauge Ψ & \texttt{psi\_for()} \\
\texttt{CSAT.md} & chaotic metastability & Lorenz, opt-in \\
\texttt{mrs.md} & bigram recognition & \texttt{bigram\_bind\_w}
(emission half retired) \\
\texttt{cll.md} & language locking & \texttt{src/lang.hpp} \\
\texttt{graph\_fisher.md} & Fisher metric &
\texttt{fisher\_distance()} \\
\texttt{RADE.md}, \texttt{maa.md}, \texttt{l9.md} & discovery, dispatch,
L9 fixes & \texttt{rade.hpp}, \texttt{mal.hpp} \\
\texttt{bmsa\_new.md}, \texttt{phenomenal-consciousness.md} & bounded
self-model & \texttt{bsma\_topo.hpp} \\
\texttt{liouville\_analog.md}, \texttt{cr\_bound.md} & proofs & no code
required \\
\end{longtable}

\subparagraph{6.2.2 Partly works --- retired as architectures, mined for
parts}\label{partly-works-retired-as-architectures-mined-for-parts}

(\texttt{research-paper/failed\_but\_parts\_works/})

These described the \textbf{DCHL substrate}, a continuous-state
Hamiltonian system that preceded the oscillator substrate. That
architecture is retired: it collapsed into the three traps of §1.2,
exactly as §2 predicts. What survived was extracted.

\begin{longtable}[]{@{}
  >{\raggedright\arraybackslash}p{(\linewidth - 4\tabcolsep) * \real{0.3333}}
  >{\raggedright\arraybackslash}p{(\linewidth - 4\tabcolsep) * \real{0.3333}}
  >{\raggedright\arraybackslash}p{(\linewidth - 4\tabcolsep) * \real{0.3333}}@{}}
\toprule\noalign{}
\begin{minipage}[b]{\linewidth}\raggedright
Paper
\end{minipage} & \begin{minipage}[b]{\linewidth}\raggedright
Status
\end{minipage} & \begin{minipage}[b]{\linewidth}\raggedright
What survived
\end{minipage} \\
\midrule\noalign{}
\endhead
\bottomrule\noalign{}
\endlastfoot
\texttt{cma.md} & superseded & manifold + Hamiltonian framing; the
vocabulary of §4 \\
\texttt{hcma.md} & superseded & DCHL contrastive training --- the
\emph{negative} result of §2.1 \\
\texttt{cch.md} & \textbf{not in live code} & curl-free constraint; grep
finds no \texttt{curl} in the substrate \\
\texttt{cmc.md}, \texttt{cmc\_extra.md} & \textbf{not in live code} &
training-memory compression; no current trainer needs it \\
\texttt{cme.md} & \textbf{never implemented} & multimodal unified state;
aspirational \\
\texttt{rifa.md} & \textbf{LIVE} (duplicated to root) & IC/SIL are
load-bearing; this copy is redundant \\
\texttt{README.md} & historical & describes the retired DCHL chatbot \\
\end{longtable}

The honest summary: this folder is the record of a substrate that did
not work and the reasons why. §2 of this paper is a self-critique built
on it. Keeping it is correct --- the negative result is a contribution
--- but nothing in it except RIFA is load-bearing.

\subparagraph{6.2.3 DEAD --- formally
retired}\label{dead-formally-retired}

(\texttt{research-paper/failed\_but\_parts\_works/DEAD/})

\textbf{These papers are dead. They are retained for provenance only and
should not be cited as describing the current system.}

\begin{longtable}[]{@{}
  >{\raggedright\arraybackslash}p{(\linewidth - 2\tabcolsep) * \real{0.5000}}
  >{\raggedright\arraybackslash}p{(\linewidth - 2\tabcolsep) * \real{0.5000}}@{}}
\toprule\noalign{}
\begin{minipage}[b]{\linewidth}\raggedright
File
\end{minipage} & \begin{minipage}[b]{\linewidth}\raggedright
Why dead
\end{minipage} \\
\midrule\noalign{}
\endhead
\bottomrule\noalign{}
\endlastfoot
\texttt{bug\_fix.md}, \texttt{future\_fix.md} & roadmaps fully
superseded by \texttt{ROADMAP.md} \\
\texttt{deprected/bed.md} & Byte Energy Descent --- converged to
byte-mean, no prompt structure (§2.1) \\
\texttt{deprected/scer.md} & multi-mode routing --- subsumed by the
field-observable dispatcher \\
\texttt{deprected/hmhra.md}, \texttt{hrm.md}, \texttt{omt.md},
\texttt{cb.md}, \texttt{cg.md} & superseded before implementation \\
\texttt{deprected/fixes.md} & ad-hoc patch log, no longer navigable \\
\end{longtable}

\subparagraph{6.2.4 Two mechanisms that outlived their
papers}\label{two-mechanisms-that-outlived-their-papers}

This is the finding the audit exists to surface.

\textbf{ZeroGrad} (\texttt{DEAD/deprected/zerograd.md}). The paper is
dead --- it described backprop-free training of the DCHL potential, and
that substrate is retired. \textbf{The learning rule is live.}
\texttt{physics.hpp} line 20 carries it verbatim:

\begin{verbatim}
ZeroGrad : plasticity_delta = pd·hebbian_decay + plastic_rate·v_a·v_b
\end{verbatim}

Every Hebbian update in the current engine, including the
phase-consistent extension of \texttt{cpme\_ii.md}, is this rule. The
paper died; the rule became foundational.

\textbf{ANQ} (\texttt{DEAD/deprected/anq.md}). The 4-bit quantisation
scheme is retired with the GPU trainer it served. \textbf{One idea
survived}: the continuous-integrator-instead-of-threshold-commit
pattern, cited in \texttt{physics.hpp} line 21 as \texttt{ANQ\ LEA}, and
now used by every saturating accumulator in the tick.

\textbf{The lesson for citation practice.} A dead paper is not a
retracted one. These two should be cited as \emph{sources of
mechanisms}, never as \emph{descriptions of the system}. Conflating the
two is how a lineage accumulates claims it can no longer defend ---
which is the failure this audit exists to prevent.

\begin{center}\rule{0.5\linewidth}{0.5pt}\end{center}

\subsubsection{7. Empirical Program}\label{empirical-program}

The architecture makes testable predictions. We enumerate them so the
empirical validation is unambiguous.

\paragraph{7.1 Substrate-Only Baselines (Falsification of §1.2
Traps)}\label{substrate-only-baselines-falsification-of-1.2-traps}

At all tested scales (\(10^4\) -- \(10^9\) parameters), substrate-only
variants (DCHL alone, BED alone, RIFA alone) should reduce to
classifier, memorizer, or retriever behaviour under any measurable
generalization test. \textbf{This is a falsification prediction of the
Translation Capacity Lemma.} Empirical evidence to date (this paper's
supporting implementation and the accompanying limitations documented in
{[}Mahir 2026e{]}) is consistent with this prediction.

\paragraph{7.2 Substrate-Level Paraphrase
Invariance}\label{substrate-level-paraphrase-invariance}

Under Hamiltonian evolution, semantically equivalent queries should
converge to states \(\mathbf{c}^*\) closer to one another than the raw
query encodings. Concretely, for a set of paraphrase clusters
\(\{P_k\}\), define

\[
r_k \equiv \frac{\text{mean intra-cluster distance after evolution}}{\text{mean intra-cluster distance before evolution}}.
\]

The claim is \(r_k < 1\) for well-trained substrates. This is \emph{the}
work only physics-based dynamics can do; if it fails, the substrate has
no non-trivial role.

\paragraph{7.3 Layer-2 Generalization to Unseen
Combinations}\label{layer-2-generalization-to-unseen-combinations}

The graph rewriter is symbolic and deterministic; it should achieve
near-100\% accuracy on syllogisms whose \emph{rules} have been supplied
but whose \emph{specific instances} have never been trained. This is the
sharpest possible test of constructive reasoning versus memorization: an
LLM at any scale below full-corpus coverage cannot achieve this, because
it has no explicit derivation mechanism.

\paragraph{7.4 OOD Honesty}\label{ood-honesty}

When queries lie outside the ontology of concepts and rules, the
pipeline must refuse --- not confabulate. The mechanism is
coherence-thresholded lookup at Layer 3, backed by IC's contradiction
reporting. The claim is that hallucination rates should be near zero on
strictly-OOD test sets, with all non-answers being explicit refusals.

\paragraph{\texorpdfstring{7.5 Scale Behavior --- \textbf{applies to the
retired DCHL substrate
only}}{7.5 Scale Behavior --- applies to the retired DCHL substrate only}}\label{scale-behavior-applies-to-the-retired-dchl-substrate-only}

\begin{quote}
This prediction is stated in units the current architecture does not
have. The oscillator substrate has \textbf{no parameters to scale}:
cavities and tunnels are created from the ontology, and their materials
are \emph{derived} from relation fan-out rather than fitted. ``1 B
parameters in 4-bit ANQ on 8 GB VRAM'' describes the retired system.

The corresponding question for the current substrate is not parameter
count but \textbf{graph scale}, and there the honest answer is a
measured limitation rather than a prediction: at 157 k concepts / 237 k
relations the per-query cost is dominated by a dense tick over all
cavities and tunnels, because Layer 1 has no sparse wavefront iteration.
Scaling behaviour is therefore currently an \emph{engineering} bound,
not an architectural claim, and should not be reported as one.
\end{quote}

\emph{(historical)} At GPU-scale substrates (\textasciitilde1B
parameters, achievable in 4-bit ANQ on 8GB VRAM), the
paraphrase-invariance ratio \(r_k\) should improve monotonically with
substrate capacity; the constructive layer's accuracy should be
independent of substrate capacity (a symbolic engine's correctness is
not a scaling law); and the translation table's necessary size \(N\) is
a property of the domain, not the substrate.

\paragraph{7.5b Empirical: The Nine-Primitive Layer 2 and a 0.5B Layer
1}\label{b-empirical-the-nine-primitive-layer-2-and-a-0.5b-layer-1}

Two empirical facts from the accompanying implementation:

\textbf{Layer 2 (constructive substrate).} All nine primitives from §4.4
fire correctly on hand-crafted premise bases (transitive, deduction,
inheritance, negation, case-analysis, induction, analogy, recursion, and
macro-rule induction). Every derived conclusion was never asserted
directly. Every derivation is traceable to its premises via the
parent-edge chain. The macro-rule induction fires on repeated derivation
streams. This validates the §3.4 composition-over-retrieval definition
on an actual buildable engine, at zero cost of ``training data'' --- the
primitives are hand-authored (a seed of \textasciitilde1 KB of code) and
their compositions are unbounded.

\textbf{Layer 1 --- HISTORICAL, on the RETIRED DCHL substrate.} The
paragraph that follows describes the \emph{previous} Layer 1, not the
current one. It is retained because the measurement was real and the
diagnosis was correct; it must not be read as evidence for the
architecture this paper now describes. The current Layer 1 is the
oscillator substrate of §5.4c, which has \textbf{no gradient training,
no parameters, and no GPU path} --- so none of these numbers apply to
it, and the ``pending fix'' below was never propagated because the
substrate it belonged to was retired instead.

\begin{quote}
\emph{(historical)} The denoising DCHL trainer, ANQ 4-bit weights, and
pinned-memory host / device pipeline complete a 4-epoch training pass on
8 000 samples from a C4 + Pile-8 corpus in \(4\) min \(43\) s wall clock
on an RTX 4060 (8 GB VRAM), using 285 MB of that VRAM for a 539 M
parameter model (D = 512, H = 1 048 576, packed in 4-bit ANQ). Denoising
\(\Delta\Phi\) is correctly signed on every epoch on which the substrate
has not yet saturated the tanh activation. \textbf{A saturation limit is
empirically visible at \textbar h\textbar{} \(\approx 0.008\)} --- the
substrate's hidden activations sit deep in tanh's linear regime ---
traceable to an aug-norm scaling of \(\sqrt{\lVert c\rVert^2 + 1}\) that
compresses pre-activations by \(\Theta(1/\sqrt{D})\) for byte-scale
inputs. The proposed fix --- replace the divisor by
\(\sqrt{(\lVert c\rVert^2 + 1)/D}\) --- was implemented in reference
code as \texttt{pre\_inv\_scale} and never propagated to the CUDA
kernel.
\end{quote}

\textbf{Why this is reported rather than deleted.} The DCHL result is
the empirical content of §2: a physics-based substrate, trained
contrastively at 0.5 B parameters, that reduced to the classifier trap
exactly as the Translation Capacity Lemma predicts. Deleting it would
remove the evidence for this paper's central self-critique. Reporting it
as though it were current would be a fabrication. It is therefore kept,
dated, and fenced.

\textbf{What the current claims actually rest on.} Layer 2's correctness
is independent of any substrate numerics: the nine primitives are exact
symbolic operations. The paper's core claims ---
composition-over-retrieval, thinker-index growth under rule induction,
verifier soundness, and the provenance chain from a derived edge back to
its premises --- rest on Layer 2 and on the oscillator Layer 1, neither
of which involves training. Every number in §7.6--§7.7 was produced on
that stack.

\paragraph{7.6 The Thinker Index and Its Growth Under Rule
Induction}\label{the-thinker-index-and-its-growth-under-rule-induction}

The thinker index \(C(\mathcal{S})/R(\mathcal{S})\) of §3.4 is directly
measurable. For any query set \(Q\), run the pipeline, count \(\ell(q)\)
for each \(q\), and compute the ratio. The claim is:

\begin{enumerate}
\def\labelenumi{\arabic{enumi}.}
\tightlist
\item
  \textbf{Baseline.} With only the initial hand-authored rules (§5.3.1),
  \(C/R > 1\) on out-of-training queries --- the value is bounded by the
  maximum derivation length allowed by the tension budget.
\item
  \textbf{Rule induction growth.} After \(M\) queries, with rule
  induction active (§5.3.3), the primitive set \(|\Pi|\) grows by
  \(O(\log M)\) (new macro-rules registered at MDL thresholds), while
  the reachable derivation space grows exponentially in \(|\Pi|\).
  \textbf{The thinker index therefore grows monotonically over the
  system's lifetime}, with no additional training data.
\item
  \textbf{Comparison to LLMs.} An LLM's thinker index is fixed at
  training time and does not grow with usage. This is \emph{the}
  separable empirical claim distinguishing SI from AI: SI systems become
  more of a thinker over time; AI systems do not.
\end{enumerate}

These are separable empirical claims, testable in a modest research
programme.

\paragraph{7.7 Recursive Bifurcation and Analogical Bridging (RB+G /
L6-A)}\label{recursive-bifurcation-and-analogical-bridging-rbg-l6-a}

The physics substrate of Layer 1 supports two structural mechanisms that
Layer 2 uses as \emph{primitives beyond the nine of §4.4}:
\textbf{Recursive Bifurcation with Gating (RB+G)} for splitting
contradictory concepts, and \textbf{L6 Analogical Bridging} for
inserting cross-domain resonance links between fissioned siblings.

\textbf{RB+G --- the mechanism.} When a concept-cavity's incoherence
metric \(I = (\Sigma^+ \cdot \Sigma^-) / (\Sigma^+ + \Sigma^-)^2\)
sustains above 0.15 with local tension \(\tau > 0.005\) for three
structural ticks, the cavity is split into two children (α excitatory, β
inhibitory) bound by a learned low-stiffness coordination spring. Four
documented failure modes of the original RB proposal {[}Mahir 2026o,
\texttt{rb.md}{]} are addressed by four mechanisms:

\begin{enumerate}
\def\labelenumi{\arabic{enumi}.}
\tightlist
\item
  \textbf{generation cap} (\(g_{\max}=1\)) --- children of a fission are
  generation-1 and ineligible for further fission, preventing cascade;
\item
  \textbf{alias ledger} --- a parallel to \texttt{cavities{[}{]}} that
  never touches \texttt{si.graph.by\_label}; child labels are synthetic
  \texttt{\textless{}parent\textgreater{}.α} / \texttt{.β};
\item
  \textbf{energy-conserving fan-out} (\(1/\sqrt{2}\)) --- a drive on a
  fissioned parent routes to both children at \(\text{amp}/\sqrt{2}\),
  so total injected energy matches a single-cavity drive at amp;
\item
  \textbf{learned binding tunnel} --- \(k_\text{init}=1.0\),
  \(p_\text{init}=0.0\), converging under Hebbian to a coordination
  stiffness rather than RB's fiat \(k=2.5\);
\item
  \textbf{audit trail} --- every fission appended to
  \texttt{data/fission\_ledger.jsonl}.
\end{enumerate}

\textbf{20-seed reproducibility} (rounds=5, seed formula
\(s = \mathtt{0xCAF17E00} \oplus (\mathtt{0x9E3779B9}\cdot(i+1))\)):

\begin{longtable}[]{@{}
  >{\raggedright\arraybackslash}p{(\linewidth - 10\tabcolsep) * \real{0.1667}}
  >{\raggedright\arraybackslash}p{(\linewidth - 10\tabcolsep) * \real{0.1667}}
  >{\raggedright\arraybackslash}p{(\linewidth - 10\tabcolsep) * \real{0.1667}}
  >{\raggedright\arraybackslash}p{(\linewidth - 10\tabcolsep) * \real{0.1667}}
  >{\raggedright\arraybackslash}p{(\linewidth - 10\tabcolsep) * \real{0.1667}}
  >{\raggedright\arraybackslash}p{(\linewidth - 10\tabcolsep) * \real{0.1667}}@{}}
\toprule\noalign{}
\begin{minipage}[b]{\linewidth}\raggedright
Concept
\end{minipage} & \begin{minipage}[b]{\linewidth}\raggedright
Cluster type
\end{minipage} & \begin{minipage}[b]{\linewidth}\raggedright
Baseline gain
\end{minipage} & \begin{minipage}[b]{\linewidth}\raggedright
RB+G gain (RMS)
\end{minipage} & \begin{minipage}[b]{\linewidth}\raggedright
Fissions/seed
\end{minipage} & \begin{minipage}[b]{\linewidth}\raggedright
Target fissioned
\end{minipage} \\
\midrule\noalign{}
\endhead
\bottomrule\noalign{}
\endlastfoot
metformin & dual-tension (act ampk, inh mtor) & 0.887 ± 0.045 & 1.517 ±
0.201 & 8.85 ± 0.96 & 20/20 \\
berberine & tri-hub (act ampk, inh mtor, inh hmg) & 0.900 ± 0.030 &
1.515 ± 0.139 & 8.60 ± 1.16 & 20/20 \\
atorvastatin & mono-directional (inh hmg, treats hcy) & 1.007 ± 0.004 &
0.957 ± 0.034 & 9.75 ± 1.84 & 20/20 \\
luna & mutual + death\_spiral inhibition & 1.244 ± 0.064 & 1.865 ± 0.162
& 7.00 ± 0.95 & 20/20 \\
linux & null (is\_a, has, no polarity conflict) & 1.032 ± 0.007 & 1.039
± 0.057 & 5.00 ± 0.00 & 20/20 \\
windows\_os & pure null (only is\_a) & 0.969 ± 0.008 & 0.955 ± 0.010 &
2.00 ± 0.00 & 0/20 \\
\end{longtable}

Concepts with polarity conflict show a large positive RB+G effect. The
mono-directional concept (atorvastatin) shows an honest null effect ---
fission fires on nearby concepts but the target's own recall does not
improve, because there was no contradiction inside it to metabolize. The
pure-null concept (windows\_os) is never fissioned in 20/20 seeds ---
the mechanism refuses to fire when the substrate has no contradiction to
resolve. This is the falsifiability guarantee: RB+G is not a free-lunch
amplifier.

\textbf{Four-lever ablation (5 seeds, metformin).} Full RB+G reaches
1.35 ± 0.05 concept RMS gain against a fission-off baseline of 0.92.

\begin{longtable}[]{@{}
  >{\raggedright\arraybackslash}p{(\linewidth - 6\tabcolsep) * \real{0.2500}}
  >{\raggedright\arraybackslash}p{(\linewidth - 6\tabcolsep) * \real{0.2500}}
  >{\raggedright\arraybackslash}p{(\linewidth - 6\tabcolsep) * \real{0.2500}}
  >{\raggedright\arraybackslash}p{(\linewidth - 6\tabcolsep) * \real{0.2500}}@{}}
\toprule\noalign{}
\begin{minipage}[b]{\linewidth}\raggedright
Configuration
\end{minipage} & \begin{minipage}[b]{\linewidth}\raggedright
RMS gain
\end{minipage} & \begin{minipage}[b]{\linewidth}\raggedright
Fissions
\end{minipage} & \begin{minipage}[b]{\linewidth}\raggedright
Reading
\end{minipage} \\
\midrule\noalign{}
\endhead
\bottomrule\noalign{}
\endlastfoot
BASELINE (fission off) & 0.922 ± 0.011 & 0.0 & negative control \\
ABLATE\_CAP (\(g_{\max}=999\)) & 1.536 ± 0.239 & 25.6 ± 3.3 & cascade
--- 3× fissions, no accuracy win \\
ABLATE\_ALIAS & 2.495 ± 0.144 & 3.4 ± 0.8 & confound: probe hits ghost
cavity \\
ABLATE\_FANOUT (factor=1.0) & 1.856 ± 0.069 & 13.6 ± 1.4 & inflated peak
--- energy doubling \\
ABLATE\_BIND & 1.470 ± 0.057 & 8.8 ± 0.8 & minor loss vs FULL \\
FULL (all four mechanisms) & 1.345 ± 0.054 & 8.4 ± 1.0 & RB+G \\
\end{longtable}

The ablation is honest, not marketing: two of the ablations produce
\emph{higher} raw RMS than FULL, and we do not hide that.
\texttt{ABLATE\_ALIAS} produces higher raw numbers because the probe
measurement itself is altered when alias routing is off --- the probe
hits the zeroed ghost cavity and the recorded post/pre ratio reflects
neighbour drift, not target recall. \texttt{ABLATE\_FANOUT} inflates the
peak because \(\sqrt{2}\times\) more energy is injected. These are
documented confounds. What FULL buys is honest energy accounting and
cascade-free convergence, not the biggest number.

\textbf{L6 Analogical Bridging.} After RB+G fissions two related
concepts, \texttt{discover\_analogies()} computes the intersection of
α-child excitatory targets and inserts a rubber tunnel (\(k=0.05\),
\(c=0.15\)) between α-cavities when \(\geq K\) targets are shared. This
is the physics realization of L6 Analogical Mapping from \texttt{l9.md}
and the structural- analogy prescription of \texttt{mims.md}. The rubber
material means MIMS impedance filtering: high-frequency mechanistic
probes are absorbed by the low-stiffness edge; low-frequency structural
probes ride the bridge.

On the insulin-pathway ontology, training
\texttt{metformin\ +\ berberine} surfaces five real analogies including:

\begin{itemize}
\tightlist
\item
  \textbf{\texttt{akt\ \textasciitilde{}\ ampk}} --- both kinases
  activating \texttt{glut4} (glucose uptake);
\item
  \textbf{\texttt{akt\ \textasciitilde{}\ mtor}} --- both kinases in the
  insulin cascade;
\item
  \textbf{\texttt{berberine\ \textasciitilde{}\ resveratrol}} --- both
  AMPK activators (a known drug-repurposing pair);
\item
  \textbf{\texttt{type2\_diabetes\ \textasciitilde{}\ drug}} (six shared
  targets) --- cross-category structural analogy;
\item
  \textbf{\texttt{molecule\ \textasciitilde{}\ ampk}} (empty shared
  target list) --- noise / weak analogy reported honestly.
\end{itemize}

Probing \texttt{metformin} at 3 Hz after bridge insertion surfaces
\texttt{atorvastatin} at rank 10 --- a concept \textbf{unreachable by
naive graph propagation} from metformin (no direct or 2-hop path in the
axioms), reachable only through the physics-inserted bridge
\texttt{metformin.α\ ↔\ berberine.α\ ↔\ hmg\_coa\_reductase\ ↔\ atorvastatin}.
In a negative-control run that fissions both concepts but skips
\texttt{discover\_analogies}, \texttt{atorvastatin} disappears from the
top-10. That single row is the L6-A signal.

\textbf{No line in any axiom file says
\texttt{metformin\ \textasciitilde{}\ berberine}}, no cavity
\texttt{metformin.α} exists at bind time, and no rubber tunnel exists
until \texttt{discover\_analogies} inserts it. This is construction over
retrieval --- the composition-over-retrieval definition of §3.4
exhibited in physics, not in symbolic proof.

\textbf{Storage-floor honesty.} The primitive facts
(\texttt{metformin\ activates\ ampk},
\texttt{berberine\ activates\ ampk}) ARE stored in
\texttt{data/insulin\_pathway.axioms}. The lemma of §3.1 requires this.
What is \emph{not} stored: the analogy
\texttt{metformin\ \textasciitilde{}\ berberine}, the child cavity
\texttt{metformin.α}, the rubber bridge, or the answer that a 3 Hz probe
at metformin surfaces atorvastatin. Every one of these is
physics-constructed from the axiom base at query time. This is what
makes it a thinker under the §3.4 definition --- a constant primitive
count, an exponentially larger reachable derivation space.

\begin{center}\rule{0.5\linewidth}{0.5pt}\end{center}

\subsubsection{8. Open Problems --- Status After RB+G /
L6-A}\label{open-problems-status-after-rbg-l6-a}

The five open problems posed in earlier drafts of this paper have been
partially resolved by the RB+G and L6-A work of §7.7. We reproduce each
open problem verbatim below and mark its current status.

\paragraph{8.1 Partially Resolved}\label{partially-resolved}

\begin{itemize}
\item
  \textbf{Domain-independent rule induction.} \emph{(Original: ``§5.3.3
  registers macro-rules from frequent derivation sub-chains. Whether
  this suffices to discover genuinely new rule families (not just macros
  of existing families) is open. Can rule induction find analogy or
  abduction rules if the initial set had only transitive
  composition?'')} \textbf{Status: RESOLVED for analogy.} L6-A (§7.7,
  \texttt{rbg.md} §5) demonstrates that the physics substrate discovers
  analogy relations \emph{without} an analogy primitive being seeded ---
  the mechanism is: RB+G fissions two concepts under contradiction;
  \texttt{discover\_analogies()} computes the intersection of their
  α-child excitatory targets; a rubber tunnel is inserted when the
  intersection is nontrivial. On the insulin-pathway ontology this
  produces \texttt{akt\ \textasciitilde{}\ ampk},
  \texttt{berberine\ \textasciitilde{}\ resveratrol}, and three other
  real analogies --- none supplied as premises. Abduction rule discovery
  from primitives-only remains open.
\item
  \textbf{Cross-domain heuristic transfer.} \emph{(Original: ``A
  substrate trained on derivation traces from one domain should provide
  some relevance signal in a related domain via shared type structure.
  Quantifying the transfer is open.'')} \textbf{Status: PARTIALLY
  RESOLVED.} L6-A demonstrates \emph{cross-concept} transfer
  \emph{within} a domain: a 3 Hz probe on \texttt{metformin} surfaces
  \texttt{atorvastatin} at rank 10 via the physics-inserted bridge
  \texttt{metformin.α\ ↔\ berberine.α\ ↔\ hmg\_coa\_reductase\ ↔\ atorvastatin}.
  Atorvastatin is unreachable from metformin by any 2-hop path in the
  raw axiom graph. This is transfer through shared \emph{structural}
  targets (both drugs act on \texttt{hmg\_coa\_reductase}), which is the
  exact ``shared type structure'' the original problem referenced. Full
  cross-\emph{domain} transfer (syllogisms → algebra) remains open and
  requires a substrate spanning multiple ontologies.
\item
  \textbf{The tension-bounded conservation law.} \emph{(Original: ``The
  expansion principle \(\Delta|G_t| \leq \alpha T(t)\) is proposed but
  not proved to be strictly conservative. What is the analog of
  Liouville's theorem for tension-bounded phase-space expansion?'')}
  \textbf{Status: EMPIRICALLY REALIZED, THEORETICALLY OPEN.} RB+G is a
  concrete mechanism whose \emph{implementation} preserves energy
  exactly: the \(1/\sqrt{2}\) fan-out (M3 of §7.7) guarantees that a
  drive on a fissioned parent injects the same total kinetic energy as a
  drive on a single cavity --- proven algebraically by
  \(2 \cdot \frac{1}{2} m (\text{amp}/\sqrt{2})^2 = \frac{1}{2} m \, \text{amp}^2\)
  and verified in the ABLATE\_FANOUT ablation (§7.7): setting fan-out to
  1.0 injects \(\sqrt{2}\times\) more energy and inflates the RMS gain
  by the same factor, confirming the fan-out is what preserves
  conservation. The generation cap \(g_\text{max} = 1\) bounds
  cumulative growth to \(\Delta|G_t| \leq 3\) per fission event (2 new
  cavities + 1 binding tunnel), against a tension budget indexed by the
  streak counter. The \emph{formal} Liouville-analog theorem --- showing
  that the tension-bounded expansion principle admits a conserved
  symplectic-like form on graph-valued phase spaces --- remains open.
\end{itemize}

\paragraph{8.2 Newly Resolved by the RB+G/L6-A
Work}\label{newly-resolved-by-the-rbgl6-a-work}

\begin{itemize}
\item
  \textbf{Contradiction metabolism as construction, not just filtering.}
  Prior to RB+G, RIFA's Integrity Core reported a scalar tension
  \(T(t)\) but did not act on it. RB+G converts sustained cavity-local
  incoherence into a \emph{physical structural mutation} --- the graph
  gains vertices and edges under contradiction rather than damping
  toward equilibrium. This is the first Layer-1 mechanism that realizes
  Axiom II operationally.
\item
  \textbf{The falsifiability guarantee for structural mechanisms.} RB+G
  on \texttt{windows\_os} (pure null control, only \texttt{is\_a} edges)
  fires zero target-fissions across 20/20 seeds. The mechanism refuses
  to fire when there is no contradiction to metabolize. This is the
  empirical form of the falsifiability criterion demanded by §7.1: a
  Layer-1 mechanism that triggers on real signal and stays silent on
  noise.
\item
  \textbf{Honest ablation methodology for physics substrates.} The
  four-lever ablation of §7.7 shows \emph{two configurations produce
  higher raw RMS than FULL} --- \texttt{ABLATE\_ALIAS} (measurement
  confound) and \texttt{ABLATE\_FANOUT} (energy inflation). Both are
  reported openly rather than hidden. This establishes a methodology for
  the empirical program of §7 in which ``the biggest number wins'' is
  not the criterion; interpretability of the measurement is.
\end{itemize}

\paragraph{8.3 Still Open}\label{still-open}

\begin{itemize}
\item
  \textbf{The Fisher-information metric on the concept graph.}
  \emph{(Original: ``CCH proposed that the cognitive manifold carries a
  Fisher-information metric. The extension to a graph-valued state
  requires the appropriate categorical generalization. Open.'')}
  \textbf{Status: STILL OPEN.} RB+G/L6-A did not address this. The
  Fisher- information metric on graph-valued states is well-defined for
  Bayesian belief networks (Amari 1985) but not for the typed rewrite
  graphs of Axiom I. This remains a categorical / differential-geometric
  question the empirical work does not touch.
\item
  \textbf{The upper bound on the thinker index.} \emph{(Original: ``Is
  \(C/R\) bounded from above for any physically realizable SI
  system?'')} \textbf{Status: STILL OPEN.} The 20-seed data of §7.7
  shows RB+G contributes a \textasciitilde1.5× RMS gain per
  contradiction-bearing concept over baseline. This does not answer the
  \emph{upper bound} question. The thinker index growth is monotone
  under rule induction (§7.6.2) and now also under L6-A bridge
  insertion, but neither implies a supremum.
\end{itemize}

\paragraph{8.4 Newly Opened by the RB+G/L6-A
Work}\label{newly-opened-by-the-rbgl6-a-work}

\begin{itemize}
\item
  \st{\textbf{Auto-tuning of fission trigger thresholds.}}
  \textbf{CLOSED.} The hand-picked values \(I > 0.15\),
  \(\tau > 0.005\), streak = 3 were replaced by homeostatic energy
  feedback: an EMA of total mechanical energy drives the threshold,
  \(I_{\text{thresh}}(t) = I_0 + k\,
  E_{\text{smooth}}(t)\). Implemented in \texttt{physics.hpp}
  (\texttt{refresh\_fission\_thresholds\_locked}). The system now sets
  its own trigger from a physical observable rather than from a constant
  tuned on one ontology.
\item
  \st{\textbf{Online L6-A discovery.}} \textbf{CLOSED.} Bridges are now
  inserted in-tick: every fission event calls
  \texttt{discover\_analogies\_locked()} restricted to the
  newly-fissioned parent, so the cost is \(O(P)\) new pairs rather than
  \(O(P^2)\) rediscovery. \texttt{layer2.hpp},
  \texttt{online\_l6a\_enabled}. The coupling with ongoing Hebbian
  consolidation that this bullet flagged as unanalysed is bounded by
  construction: bridges are rubber tunnels with \(k = 0.05\), an order
  of magnitude below the stiffest structural material, so an inserted
  bridge cannot dominate consolidation on the tick it appears.
\item
  \textbf{Rubber-tunnel frequency selectivity.} The claim that
  low-frequency probes ride the L6-A bridge while high-frequency probes
  are absorbed depends on the impedance mismatch between rubber
  (\(k=0.05\)) and steel (\(k \geq 1.6\)) tunnels. The optimal probe
  frequency 3 Hz was chosen empirically on the insulin ontology. A
  theoretical treatment (transfer- function analysis of the MIMS
  material table) would let us pick probe frequencies for arbitrary
  domains without inspection.
\item
  \textbf{Scaling to \(10^4\)--\(10^5\) concepts.} \emph{(Partly closed,
  and the remainder is now measured rather than anticipated.)} The
  \textbf{Layer-2} structural tick is sparse:
  \texttt{quiet\_ticks{[}i{]}} tracks the wavefront and the fission
  trigger evaluates only recently-driven cavities, bit-identical to the
  dense baseline over 5 seeds. \textbf{Layer 1 is not.}
  \texttt{PhysicsEngine::tick()} still iterates every cavity and every
  tunnel, so at 157 k concepts / 237 k relations a single query costs
  \textasciitilde200 ticks × 237 k tunnels × 4 transcendentals. This is
  the binding constraint on ontology size today, and it is an
  implementation gap, not an architectural one --- the wavefront data
  structure already exists and Layer 2 already honours it.

  Three \(O(E^2)\) defects on the ingest path were found and fixed while
  measuring this: \texttt{has\_edge} was a linear scan called from every
  \texttt{add\_edge} (making a bulk load quadratic),
  \texttt{would\_contradict} scanned all edges per asserted axiom, and
  four reasoning primitives performed their joins by full edge scan.
  Indexing them reduced load time by \(105\times\) at fixed output. None
  of this changes a claim in this paper; it changes what is runnable.
\item
  \st{\textbf{Language generation Stage 2 (syntax macro-layer).}}
  \textbf{ABANDONED, not pending.} This bullet described Stage 2 of
  §5.4b as ``a small extension not yet implemented''. It was in fact
  implemented and it failed: the measured slot readout left the Subject
  or Verb slot \emph{empty} on 3 of 3 probes. More importantly the
  approach was \textbf{retired as a category error} (§5.4d): grammatical
  role is a convention, and a fixed natural frequency \(\omega_S\)
  cannot discover a convention. The question this bullet posed --- ``is
  English syntax captured by three-slot phase-locking?'' --- is not
  open. It is answered \emph{no}, and the replacement is a deterministic
  protocol table (\texttt{cll.md}).
\end{itemize}

\begin{center}\rule{0.5\linewidth}{0.5pt}\end{center}

\paragraph{8.5 The Central Unsolved Problem: Continuous Wave →
Syntax-Constrained Byte
Stream}\label{the-central-unsolved-problem-continuous-wave-syntax-constrained-byte-stream}

Everything above concerns \emph{what the field knows} and \emph{how it
changes}. This section concerns \emph{how the field speaks}, and it is
the one place where we have a mechanism specification, a measured
failure, and no solution.

\textbf{The problem.} Layer 1 holds meaning as a standing-wave field.
Layer 3 converts a graph invariant to a stored string. Between them sits
a gap: how does field energy become an \textbf{ordered byte stream that
is grammatical}, without collapsing into babble?

\[
\text{Wave dynamics in cavity field}
\;\xrightarrow{\;\text{phase locking + gauge coupling}\;}\;
\text{syntax-constrained mechanical emission}
\]

\textbf{Why the field ran to language models instead.} This is worth
stating plainly, because it explains the shape of the entire discipline.
The statistical shortcut

\[
P(w_n \mid w_1 \ldots w_{n-1})
\]

is not a theory of language production. It is what remains once the
mechanical problem is abandoned. A next-token distribution requires no
account of \emph{why} a word follows another --- only a count of how
often it did. Every architecture that adopts it has, at that moment,
stopped attempting emission physics and started interpolating a corpus.
The approach works, spectacularly, and it forecloses the question.

The claim of this paper is not that the statistical route is wrong. It
is that the mechanical route was never attempted seriously, and that its
difficulty --- not its impossibility --- is why.

\textbf{What is built.} Word \emph{order across concepts} is solved.
CPME §1.6 sweeps a carrier phase across a per-cavity phase field
\(\phi_i = \operatorname{atan2}(v_i, \omega_0 x_i)\); CPSB §2.1 supplies
direction through per-tunnel gauge offsets \(\Psi\); output length is
emergent rather than templated. On code structure this reaches a correct
signature reconstruction that compiles (§5.4c). At the level of
\emph{which concept follows which}, physics-driven ordering works.

\textbf{What is not built.} Word \emph{form} --- the byte stream itself.
MRS §1.4 proposes emitting bytes by sweeping the byte-lane phase field.
It babbles. The measured failures are recorded in \texttt{mrs.md} §4b,
and the recognition half of that paper (bigram lanes, 7/7 on
misspellings) must not be mistaken for the emission half, which does not
work.

\textbf{Why the cheap fixes are known not to work.} Two were tried and
both failed for the same structural reason:

\begin{itemize}
\tightlist
\item
  \emph{Reweighting the byte scoring.} Every reweighting relocated the
  failure rather than removing it --- the signature of insufficient
  representation, not of bad parameters.
\item
  \emph{Widening the alphabet to bigrams.} This fixed
  \textbf{recognition}, because recognition compares a query against
  stored forms. It does not fix \textbf{emission}, because emission has
  nothing to compare against: the field must \emph{originate} a
  sequence, not match one.
\end{itemize}

The asymmetry is the point. Recognition is a comparison problem and
yields to better representation. Emission is a \textbf{generation}
problem and requires a mechanism that converts continuous field energy
into discrete ordered symbols under grammatical constraint. No amount of
representational widening supplies that mechanism.

\textbf{REVISION (§5.4d): the problem was two problems, and one is
solved.}

This section originally asked for grammatical constraint as a
\emph{field property}. That was a category error. Grammar is a
transmission protocol, not a phenomenon, and §5.4d separates the two:

\begin{longtable}[]{@{}
  >{\raggedright\arraybackslash}p{(\linewidth - 4\tabcolsep) * \real{0.3333}}
  >{\raggedright\arraybackslash}p{(\linewidth - 4\tabcolsep) * \real{0.3333}}
  >{\raggedright\arraybackslash}p{(\linewidth - 4\tabcolsep) * \real{0.3333}}@{}}
\toprule\noalign{}
\begin{minipage}[b]{\linewidth}\raggedright
\end{minipage} & \begin{minipage}[b]{\linewidth}\raggedright
mechanism
\end{minipage} & \begin{minipage}[b]{\linewidth}\raggedright
status
\end{minipage} \\
\midrule\noalign{}
\endhead
\bottomrule\noalign{}
\endlastfoot
\textbf{Which concepts, in what relation} & physics --- resonance, gauge
Ψ, phase-collapse & solved \\
\textbf{Word order, article, script} & protocol --- isolated language
cluster & \textbf{solved (§5.4d)} \\
\textbf{Word FORM for a lexicon-absent word} & requires emission
mechanics & \textbf{open} \\
\end{longtable}

Ordering across concepts and realisation into English, Bengali or
Banglish are done, deterministically and without ML. What remains is
strictly narrower than originally stated: producing the \emph{bytes of a
word the lexicon does not contain}. For any word already in Layer 3, no
byte synthesis is needed --- the protocol layer emits it exactly.

The residual problem therefore matters for morphological productivity
(inflecting an unseen stem, coining a compound) rather than for speaking
at all. That is a real limitation, and a much smaller one.

\textbf{What the residual problem actually requires.} Not a trick --- an
emission mechanics for the field itself. Minimally:

\begin{enumerate}
\def\labelenumi{\arabic{enumi}.}
\tightlist
\item
  \textbf{A discretisation principle.} What physical event constitutes
  ``emit a byte now''? Threshold crossing, phase coincidence, and
  resonance capture are candidates; none has been derived.
\item
  \textbf{A refractory / sequencing law.} What prevents the same lane
  from dominating, and what advances the field to the next symbol? The
  current sweep imposes this externally rather than deriving it.
\item
  \textbf{A grammatical constraint expressed as a field property.} CPSB
  shows this is possible at the concept level (\(\Psi\) as a directed
  phase offset). The morphological analogue --- a constraint that makes
  ill-formed byte sequences physically unfavourable rather than merely
  unlikely --- does not exist.
\item
  \textbf{A termination condition.} When is an utterance complete?
  Silence must be a field state, not a counter reaching zero.
\end{enumerate}

\textbf{Status.} Open, and correctly identified as the hard core rather
than a remaining detail. It is a mathematical and physical problem, not
an engineering one, and it should not be reported as solved until a
mechanism satisfying (1)--(4) is specified and measured. The honest
position is that this substrate currently \emph{thinks} in fields and
\emph{speaks} through a stored label table --- and that the second half
is exactly what remains.

\begin{center}\rule{0.5\linewidth}{0.5pt}\end{center}

\subsubsection{9. Conclusion}\label{conclusion}

Synthetic Intelligence is not scaled Artificial Intelligence. SI is a
\emph{categorically different mode of computation}: active,
constructive, structurally grounded, and information-theoretically
honest about where its knowledge lives. The Translation Capacity Lemma
shows that no cognitive architecture can escape the mandatory
\(\log_2 N\) cost of discriminating \(N\) semantic outcomes; the honest
response is not to hide this cost but to name it, minimize it, and pair
it with a substrate that does what only physics-based dynamics can do.

The Tri-Layer Unified Architecture is the concrete realization: a
physics-based coherence substrate that no lookup table can replace, a
graph-rewriting constructive layer that no continuous dynamics can
perform, and a minimum-description-length translation table that no
distributed embedding can match on interpretability. Each layer is
minimal. Each layer has a specific, testable role. Removing any layer
costs a specific capability. The synthesis of our prior lineage into
this architecture is the argument of this paper.

The path from AI to SI is not through more scale. It is through the
recognition that thinking is construction, not selection; that
intelligence is structural synthesis, not pattern matching; and that any
architecture that pretends to eliminate the translation layer is simply
hiding it. The honest architecture names its costs, distributes its
responsibilities, and builds cognition where cognition can be built ---
not where it can only be simulated.

\begin{center}\rule{0.5\linewidth}{0.5pt}\end{center}

\subsubsection{References}\label{references}

{[}Mahir 2026a{]} The Cognitive Manifold Architecture.
\emph{research-paper/cma.md}. {[}Mahir 2026b{]} Predictive Cognitive
Manifold Architecture (P-CMA). \emph{research-paper/pcma.md}. {[}Mahir
2026c{]} The Cognitive Continuum Hypothesis (CCH).
\emph{research-paper/cch.md}. {[}Mahir 2026d{]} The ZeroGrad Cognitive
Engine. \emph{research-paper/zerograd.md}. {[}Mahir 2026e{]} Byte Energy
Descent (BED). \emph{research-paper/bed.md}. {[}Mahir 2026f{]}
Arch-Native Quantization (ANQ). \emph{research-paper/anq.md}. {[}Mahir
2026g{]} Recursive Integrity Field Architecture (RIFA).
\emph{research-paper/deprected/rifa.md}. {[}Mahir 2026h{]} Hierarchical
Cognitive Manifold Architecture (HCMA). \emph{research-paper/hcma.md}.
{[}Mahir 2026i{]} The Cognitive Manifold Compressor (CMC).
\emph{research-paper/cmc.md}. {[}Mahir 2026j{]} The Cognitive Multimodal
Engine (CME). \emph{research-paper/cme.md}. {[}Mahir 2026k{]} LTO: Live
Training Over Manifolds. \emph{research-paper/lto.md}. {[}Mahir 2026l{]}
SCER: Search + Compute + Evolve + Reason. \emph{research-paper/scer.md}.
{[}Besag 1986{]} On the Statistical Analysis of Dirty Pictures. \emph{J.
R. Statist. Soc. B}. {[}Lawvere 1963{]} Functorial Semantics of
Algebraic Theories. \emph{PhD thesis, Columbia}. {[}Mac Lane 1971{]}
Categories for the Working Mathematician. \emph{Springer}. {[}Shannon
1948{]} A Mathematical Theory of Communication. \emph{Bell Sys. Tech.
J.} {[}Kolmogorov 1965{]} Three Approaches to the Quantitative
Definition of Information. {[}Grothendieck 1957{]} Sur quelques points
d'algèbre homologique. \emph{Tohoku Math. J.} {[}LeCun 2022{]} A Path
Towards Autonomous Machine Intelligence. {[}Marr 1982{]} Vision: A
Computational Investigation. \emph{W. H. Freeman}. {[}Newell \& Simon
1976{]} Computer Science as Empirical Inquiry. \emph{Comm. ACM}.
{[}Turing 1950{]} Computing Machinery and Intelligence. \emph{Mind}.

\begin{center}\rule{0.5\linewidth}{0.5pt}\end{center}

\textbf{From Artificial to Synthetic Intelligence: pattern matching is
not thinking, retrieval is not reasoning, and the honest cognitive
architecture is the one that names its costs, pays them minimally, and
builds construction where construction can be built.}

\end{document}
